% Generated by tools/build_appendix.py; do not hand-edit.
% Styling is owned exclusively by 章节模板/附录模板.tex.
\AppendixStart
\AppendixPart{1}{附件文件及内容}
\AppendixAttachmentTableBegin
\AppendixAttachmentRow{1}{问题一}{完整方法源码、代码环境说明和 Task/调度逐例结果表。}
\AppendixAttachmentRow{2}{问题二}{完整方法源码、代码环境说明和跨核交接/DDR 逐例结果表。}
\AppendixAttachmentRow{3}{问题三}{完整方法源码、代码环境说明和 L2/FIFO 同方案配对逐例结果表。}
\AppendixAttachmentRow{4}{其他数据}{100 个原始计算图、官方配置、完整官方评价代码及来源说明。}
\AppendixAttachmentRow{5}{其他支撑材料}{1400 个已选主配置和 500 个同方案配对的计划/官方响应、各分片执行源码、数据表及整理复现程序。}
\AppendixAttachmentTableEnd
\AppendixPart{2}{问题一关键代码}
\AppendixCodeFile[firstline=58,lastline=121,firstnumber=58]{官方评价与字典序接受（片段 1）}{python}{../提交附件/支撑材料文件/问题一/源码/research_a_orion_trident/solve.py}
\AppendixCodeFile[firstline=43,lastline=49,firstnumber=43]{初始结构候选构造（片段 2）}{python}{../提交附件/支撑材料文件/问题一/源码/research_a_orion_trident/sources.py}
\AppendixPart{3}{问题二关键代码}
\AppendixCodeFile[firstline=350,lastline=376,firstnumber=350]{跨核交接与 DDR 搬运量（片段 1）}{python}{../提交附件/支撑材料文件/问题二/源码/research_a_orion_trident/vendor/orion/core.py}
\AppendixCodeFile[firstline=377,lastline=431,firstnumber=377]{全局及候选下界（片段 2）}{python}{../提交附件/支撑材料文件/问题二/源码/research_a_orion_trident/vendor/orion/core.py}
\AppendixPart{4}{问题三关键代码}
\AppendixCodeFile[firstline=19,lastline=63,firstnumber=19]{FIFO 事件驱动的候选定位与顺序移动（片段 1）}{python}{../提交附件/支撑材料文件/问题三/源码/research_a_orion_trident/fifo.py}
\AppendixCodeFile[firstline=258,lastline=275,firstnumber=258]{同一完整方案的无 L2 官方配对（片段 2）}{python}{../提交附件/支撑材料文件/问题三/源码/research_a_orion_trident/solve.py}
\AppendixPart{5}{逐例结果与同方案配对}
% 可选附录长表；每表最多两个核数，表宽不超过 160 mm。

% 问题一、二的 1 核行从规范化 appendix_q1/q2.csv 中题设单核约定读取，不由多核公式回填。

% 周期为总体任务执行时间；字节为额外 COPY。q1/q2 单核使用题设约定的原始整图响应。
\begingroup\normalsize
\setlength{\tabcolsep}{3pt}
\setlength{\LTleft}{\fill}\setlength{\LTright}{\fill}
\setlength{\LTcapwidth}{160mm}
\begin{longtable}{@{}p{15mm}>{\raggedleft\arraybackslash}p{28mm}>{\raggedleft\arraybackslash}p{39mm}>{\raggedleft\arraybackslash}p{28mm}>{\raggedleft\arraybackslash}p{39mm}@{}}
\caption{问题一逐例结果：1、2 核}\label{tab:appendix_q1_n12}\\
\toprule
算例 & \shortstack{1 核\\周期} & \shortstack{1 核额外\\COPY/B} & \shortstack{2 核\\周期} & \shortstack{2 核额外\\COPY/B}\\
\midrule
\endfirsthead
\caption[]{问题一逐例结果：1、2 核（续）}\\
\toprule
算例 & \shortstack{1 核\\周期} & \shortstack{1 核额外\\COPY/B} & \shortstack{2 核\\周期} & \shortstack{2 核额外\\COPY/B}\\
\midrule
\endhead
\bottomrule
\endfoot
001 & 233,110 & 0 & 116,868 & 1,152\\
002 & 261,945 & 0 & 130,549 & 1,585,152\\
003 & 602,212 & 0 & 356,773 & 5,569,332\\
004 & 392,357 & 0 & 199,074 & 110,592\\
005 & 95,618 & 0 & 62,870 & 1,034,798\\
006 & 82,837 & 16,400 & 44,908 & 24,448\\
007 & 1,012,562 & 0 & 507,494 & 27,648\\
008 & 487,605 & 0 & 137,478 & 3,981,312\\
009 & 153,261 & 0 & 91,786 & 214,756\\
010 & 86,969 & 0 & 41,864 & 1,151,576\\
011 & 184,697 & 0 & 92,646 & 180,224\\
012 & 59,088 & 0 & 29,814 & 250,046\\
013 & 263,949 & 0 & 132,198 & 8,192\\
014 & 17,698,626 & 71,010,720 & 9,159,748 & 165,905,592\\
015 & 177,362 & 180,256 & 91,347 & 240,920\\
016 & 7,715,523 & 0 & 4,550,194 & 119,554,128\\
017 & 162,810 & 22,528 & 80,376 & 1,581,434\\
018 & 1,004,367 & 1,970,176 & 507,027 & 2,082,560\\
019 & 81,144 & 0 & 41,573 & 757,344\\
020 & 532,746 & 0 & 266,580 & 0\\
021 & 6,567,937 & 4,501,504 & 3,150,849 & 18,173,952\\
022 & 56,780 & 0 & 28,304 & 6,912\\
023 & 122,614 & 0 & 53,954 & 1,660,704\\
024 & 2,555,343 & 0 & 1,506,718 & 39,327,048\\
025 & 4,863,026 & 0 & 2,431,840 & 27,648\\
026 & 99,203 & 0 & 51,136 & 19,488\\
027 & 2,363,952 & 1,400,832 & 1,120,962 & 1,744,896\\
028 & 43,836,040 & 277,412,256 & 20,164,856 & 122,671,488\\
029 & 37,555 & 0 & 19,740 & 0\\
030 & 2,808,034 & 5,519,360 & 1,418,463 & 3,617,492\\
031 & 1,414,369 & 3,491,328 & 711,541 & 6,832,674\\
032 & 41,443 & 0 & 21,170 & 6,216\\
033 & 1,774,196 & 4,592,640 & 893,967 & 6,776,976\\
034 & 842,996 & 1,079,832 & 423,014 & 2,613,828\\
035 & 169,571 & 0 & 95,406 & 0\\
036 & 931,510 & 0 & 466,068 & 1,152\\
037 & 210,001 & 0 & 105,374 & 204,800\\
038 & 1,364,555 & 5,627,904 & 677,057 & 6,666,248\\
039 & 1,435,262 & 7,929,856 & 711,843 & 28,516,352\\
040 & 253,454 & 0 & 144,951 & 12,288\\
041 & 6,003,957 & 30,898,176 & 3,036,273 & 19,746,124\\
042 & 1,686,740 & 6,629,376 & 853,276 & 16,912,800\\
043 & 677,656 & 49,160 & 342,474 & 1,951,562\\
044 & 154,407 & 4,438,112 & 68,478 & 1,410,176\\
045 & 125,614 & 0 & 63,059 & 0\\
046 & 276,455 & 3,176,288 & 124,461 & 2,152,480\\
047 & 326,508 & 0 & 206,942 & 2,585,394\\
048 & 222,220 & 0 & 165,474 & 876,312\\
049 & 196,276 & 0 & 129,592 & 1,105,746\\
050 & 149,690 & 0 & 94,785 & 948,954\\
051 & 607,628 & 0 & 357,955 & 9,045,258\\
052 & 180,301 & 0 & 89,547 & 811,806\\
053 & 1,472,686 & 5,405,536 & 755,778 & 13,628,414\\
054 & 815,051 & 0 & 456,207 & 3,233,186\\
055 & 118,835 & 69,632 & 59,573 & 57,856\\
056 & 277,285 & 0 & 199,146 & 2,101,798\\
057 & 59,726 & 94,216 & 31,741 & 80,512\\
058 & 4,539,717 & 28,696,576 & 2,240,589 & 96,509,952\\
059 & 829,473 & 1,691,648 & 369,318 & 360,448\\
060 & 888,789 & 2,068,520 & 449,088 & 3,955,756\\
061 & 77,191 & 0 & 38,453 & 34,560\\
062 & 2,854,809 & 0 & 1,477,150 & 24,446,976\\
063 & 1,070,397 & 0 & 545,710 & 23,526,912\\
064 & 19,116 & 0 & 15,360 & 47,664\\
065 & 53,068 & 0 & 24,248 & 269,548\\
066 & 359,070 & 0 & 203,737 & 1,179,430\\
067 & 63,834,834 & 264,449,408 & 30,245,653 & 75,973,504\\
068 & 362,095 & 0 & 249,014 & 3,617,792\\
069 & 25,144 & 0 & 14,586 & 166,124\\
070 & 117,300 & 0 & 58,199 & 448,662\\
071 & 18,919 & 0 & 12,103 & 137,002\\
072 & 23,897,828 & 108,613,680 & 11,770,887 & 175,833,776\\
073 & 10,081,626 & 141,570,432 & 4,265,708 & 13,606,400\\
074 & 1,682,239 & 6,070,280 & 847,649 & 7,897,104\\
075 & 1,441,809 & 0 & 779,051 & 5,484,858\\
076 & 6,185,113 & 28,280,832 & 3,220,217 & 74,277,888\\
077 & 495,320 & 0 & 250,014 & 928,058\\
078 & 2,296,143 & 12,644,568 & 1,080,972 & 15,039,360\\
079 & 31,052,251 & 32,448,512 & 14,843,277 & 91,557,888\\
080 & 292,138 & 548,864 & 145,926 & 2,191,360\\
081 & 151,795 & 0 & 76,860 & 0\\
082 & 721,041 & 0 & 379,264 & 3,054,660\\
083 & 1,057,993 & 13,934,688 & 467,992 & 5,000,000\\
084 & 2,507,412 & 0 & 687,332 & 9,925,632\\
085 & 3,187,899 & 0 & 1,683,359 & 14,352,330\\
086 & 118,403 & 0 & 78,444 & 1,471,960\\
087 & 3,974,071 & 7,902,720 & 1,991,892 & 10,987,328\\
088 & 243,623 & 0 & 149,950 & 1,449,994\\
089 & 693,092 & 0 & 346,832 & 238,720\\
090 & 1,178,363 & 21,233,592 & 494,203 & 3,534,720\\
091 & 9,186,769 & 53,947,008 & 4,641,137 & 111,883,264\\
092 & 4,815,554 & 66,000,096 & 2,054,026 & 9,710,496\\
093 & 74,205 & 0 & 37,827 & 8,192\\
094 & 94,892 & 0 & 47,990 & 106,272\\
095 & 2,084,853 & 0 & 575,022 & 17,031,168\\
096 & 84,194 & 0 & 43,430 & 67,392\\
097 & 11,491,509 & 10,584,064 & 5,192,028 & 29,818,880\\
098 & 177,430 & 0 & 89,756 & 0\\*
099 & 1,622,999 & 5,484,544 & 809,832 & 7,330,944\\*
100 & 151,223 & 0 & 76,840 & 55,296\\
\end{longtable}
\endgroup

% 周期为总体任务执行时间；字节为额外 COPY。q1/q2 单核使用题设约定的原始整图响应。
\begingroup\normalsize
\setlength{\tabcolsep}{3pt}
\setlength{\LTleft}{\fill}\setlength{\LTright}{\fill}
\setlength{\LTcapwidth}{160mm}
\begin{longtable}{@{}p{15mm}>{\raggedleft\arraybackslash}p{28mm}>{\raggedleft\arraybackslash}p{39mm}>{\raggedleft\arraybackslash}p{28mm}>{\raggedleft\arraybackslash}p{39mm}@{}}
\caption{问题一逐例结果：3、4 核}\label{tab:appendix_q1_n34}\\
\toprule
算例 & \shortstack{3 核\\周期} & \shortstack{3 核额外\\COPY/B} & \shortstack{4 核\\周期} & \shortstack{4 核额外\\COPY/B}\\
\midrule
\endfirsthead
\caption[]{问题一逐例结果：3、4 核（续）}\\
\toprule
算例 & \shortstack{3 核\\周期} & \shortstack{3 核额外\\COPY/B} & \shortstack{4 核\\周期} & \shortstack{4 核额外\\COPY/B}\\
\midrule
\endhead
\bottomrule
\endfoot
001 & 78,545 & 2,304 & 58,984 & 3,456\\
002 & 89,993 & 1,585,152 & 70,039 & 56,832\\
003 & 304,000 & 6,453,800 & 259,077 & 6,655,708\\
004 & 134,043 & 221,184 & 104,188 & 331,776\\
005 & 55,087 & 1,071,952 & 51,058 & 1,298,726\\
006 & 31,152 & 17,280 & 24,739 & 201,000\\
007 & 338,094 & 55,296 & 254,398 & 82,944\\
008 & 113,080 & 3,981,312 & 112,880 & 3,981,312\\
009 & 91,619 & 340,494 & 89,993 & 1,198,942\\
010 & 30,925 & 345,620 & 27,887 & 527,542\\
011 & 69,625 & 1,081,344 & 46,880 & 540,672\\
012 & 20,673 & 20,160 & 16,477 & 48,384\\
013 & 88,351 & 16,384 & 66,982 & 24,576\\
014 & 6,197,366 & 165,905,592 & 4,704,368 & 165,905,592\\
015 & 64,732 & 290,056 & 45,942 & 358,920\\
016 & 3,712,678 & 119,554,736 & 3,241,338 & 119,555,344\\
017 & 54,299 & 549,320 & 42,259 & 252,288\\
018 & 344,212 & 2,954,368 & 255,738 & 3,998,848\\
019 & 28,530 & 29,760 & 21,453 & 85,060\\
020 & 177,950 & 0 & 133,704 & 0\\
021 & 2,117,298 & 18,223,104 & 1,550,324 & 4,608,000\\
022 & 19,966 & 6,912 & 15,412 & 6,912\\
023 & 43,459 & 158,112 & 31,369 & 232,128\\
024 & 1,228,774 & 39,327,248 & 1,072,818 & 39,327,448\\
025 & 1,622,164 & 55,296 & 1,217,444 & 82,944\\
026 & 40,728 & 443,132 & 37,256 & 835,222\\
027 & 824,914 & 12,691,456 & 559,196 & 3,317,760\\
028 & 13,429,199 & 123,848,568 & 10,326,529 & 124,363,008\\
029 & 12,743 & 0 & 11,028 & 0\\
030 & 949,141 & 5,453,524 & 711,582 & 7,335,636\\
031 & 492,843 & 13,306,368 & 371,121 & 8,212,920\\
032 & 15,406 & 12,432 & 12,462 & 12,360\\
033 & 595,340 & 7,480,384 & 448,056 & 7,843,424\\
034 & 290,064 & 2,765,824 & 210,056 & 3,017,216\\
035 & 95,341 & 34,068 & 95,341 & 34,068\\
036 & 311,345 & 2,304 & 233,584 & 3,456\\
037 & 74,038 & 1,064,960 & 53,320 & 614,400\\
038 & 459,698 & 8,326,664 & 347,714 & 7,068,172\\
039 & 532,103 & 19,410,944 & 403,350 & 7,725,056\\
040 & 110,105 & 0 & 110,105 & 0\\
041 & 2,001,932 & 21,414,220 & 1,553,352 & 64,857,104\\
042 & 576,279 & 17,044,896 & 436,410 & 17,044,896\\
043 & 264,649 & 2,600,012 & 233,346 & 2,204,108\\
044 & 61,498 & 1,513,376 & 59,360 & 1,619,936\\
045 & 42,601 & 0 & 31,878 & 0\\
046 & 105,123 & 3,055,712 & 95,844 & 3,837,312\\
047 & 175,702 & 2,779,246 & 162,540 & 2,955,232\\
048 & 151,336 & 925,804 & 146,785 & 947,696\\
049 & 106,320 & 1,189,054 & 95,413 & 1,917,854\\
050 & 75,063 & 727,418 & 66,866 & 1,476,570\\
051 & 291,222 & 9,045,304 & 254,308 & 9,045,350\\
052 & 61,057 & 176,640 & 47,608 & 250,368\\
053 & 552,794 & 10,563,176 & 481,965 & 12,029,352\\
054 & 380,010 & 6,147,860 & 335,487 & 6,599,942\\
055 & 42,937 & 99,328 & 33,024 & 292,868\\
056 & 160,283 & 2,107,946 & 143,316 & 3,919,844\\
057 & 23,178 & 28,736 & 16,189 & 72,832\\
058 & 1,661,359 & 39,198,720 & 1,233,005 & 14,422,016\\
059 & 260,453 & 7,743,488 & 185,216 & 1,081,344\\
060 & 307,811 & 8,098,700 & 219,705 & 3,604,480\\
061 & 27,892 & 33,792 & 21,814 & 110,336\\
062 & 966,144 & 130,560 & 762,127 & 130,560\\
063 & 378,900 & 99,840 & 275,527 & 99,840\\
064 & 13,097 & 93,012 & 12,254 & 101,978\\
065 & 18,576 & 27,648 & 14,545 & 41,616\\
066 & 183,820 & 1,311,946 & 157,421 & 3,899,786\\
067 & 20,503,402 & 86,178,560 & 15,573,083 & 82,792,576\\
068 & 203,282 & 4,543,290 & 182,864 & 4,307,652\\
069 & 10,847 & 213,840 & 10,732 & 261,940\\
070 & 41,459 & 287,752 & 31,750 & 346,624\\
071 & 9,136 & 212,698 & 8,565 & 226,228\\
072 & 7,900,998 & 175,833,776 & 5,950,651 & 125,592,070\\
073 & 4,187,840 & 15,118,976 & 2,861,738 & 17,675,392\\
074 & 564,254 & 14,401,540 & 430,147 & 9,211,920\\
075 & 670,128 & 6,039,214 & 526,344 & 6,839,538\\
076 & 2,150,470 & 74,277,888 & 1,617,717 & 74,277,888\\
077 & 178,900 & 935,736 & 134,597 & 1,015,510\\
078 & 741,065 & 15,572,832 & 619,540 & 15,039,360\\
079 & 9,899,446 & 91,557,888 & 7,454,944 & 91,680,768\\
080 & 100,325 & 3,047,424 & 100,325 & 3,047,424\\
081 & 51,943 & 0 & 39,588 & 0\\
082 & 294,320 & 3,264,534 & 270,415 & 3,667,456\\
083 & 339,303 & 4,421,216 & 259,314 & 4,495,168\\
084 & 458,734 & 9,925,632 & 392,976 & 9,925,632\\
085 & 1,302,719 & 15,781,494 & 1,137,740 & 16,158,338\\
086 & 70,592 & 1,260,420 & 62,644 & 1,403,476\\
087 & 1,378,050 & 13,704,154 & 1,081,918 & 25,367,040\\
088 & 133,515 & 2,728,012 & 117,517 & 2,279,726\\
089 & 243,614 & 477,440 & 178,160 & 716,160\\
090 & 462,845 & 3,988,896 & 396,616 & 6,970,704\\
091 & 3,106,885 & 111,883,264 & 2,359,089 & 111,883,264\\
092 & 1,412,994 & 9,069,632 & 1,085,306 & 9,479,296\\
093 & 25,103 & 16,384 & 19,779 & 24,576\\
094 & 31,921 & 211,968 & 24,974 & 318,816\\
095 & 477,352 & 17,031,168 & 477,352 & 17,031,168\\
096 & 29,834 & 67,392 & 26,861 & 661,824\\
097 & 3,532,025 & 31,096,832 & 2,608,028 & 29,818,880\\
098 & 69,230 & 1,808,822 & 59,379 & 1,792,440\\*
099 & 543,319 & 8,883,328 & 405,202 & 9,845,888\\*
100 & 58,189 & 412,176 & 57,842 & 420,390\\
\end{longtable}
\endgroup

% 周期为总体任务执行时间；字节为额外 COPY。q1/q2 单核使用题设约定的原始整图响应。
\begingroup\normalsize
\setlength{\tabcolsep}{3pt}
\setlength{\LTleft}{\fill}\setlength{\LTright}{\fill}
\setlength{\LTcapwidth}{160mm}
\begin{longtable}{@{}p{15mm}>{\raggedleft\arraybackslash}p{28mm}>{\raggedleft\arraybackslash}p{39mm}@{}}
\caption{问题一逐例结果：5 核}\label{tab:appendix_q1_n5}\\
\toprule
算例 & \shortstack{5 核\\周期} & \shortstack{5 核额外\\COPY/B}\\
\midrule
\endfirsthead
\caption[]{问题一逐例结果：5 核（续）}\\
\toprule
算例 & \shortstack{5 核\\周期} & \shortstack{5 核额外\\COPY/B}\\
\midrule
\endhead
\bottomrule
\endfoot
001 & 47,502 & 4,608\\
002 & 62,576 & 615,936\\
003 & 246,590 & 7,109,294\\
004 & 88,341 & 442,368\\
005 & 51,298 & 1,246,498\\
006 & 21,084 & 306,882\\
007 & 203,834 & 110,592\\
008 & 100,603 & 0\\
009 & 82,966 & 1,646,012\\
010 & 26,184 & 539,064\\
011 & 46,880 & 540,672\\
012 & 13,786 & 50,112\\
013 & 54,363 & 32,768\\
014 & 3,842,702 & 165,905,592\\
015 & 40,072 & 432,640\\
016 & 3,241,338 & 119,555,344\\
017 & 34,502 & 339,840\\
018 & 204,766 & 5,320,704\\
019 & 17,766 & 63,076\\
020 & 107,650 & 0\\
021 & 1,283,116 & 18,124,800\\
022 & 12,969 & 6,912\\
023 & 26,451 & 355,008\\
024 & 1,067,515 & 39,327,648\\
025 & 975,282 & 110,592\\
026 & 37,256 & 835,222\\
027 & 528,056 & 14,319,616\\
028 & 8,396,783 & 126,483,456\\
029 & 9,383 & 0\\
030 & 573,711 & 9,119,444\\
031 & 310,507 & 8,674,748\\
032 & 9,824 & 49,248\\
033 & 359,312 & 8,437,888\\
034 & 177,773 & 3,653,632\\
035 & 92,432 & 267,396\\
036 & 187,182 & 4,608\\
037 & 43,800 & 1,925,120\\
038 & 287,328 & 8,299,532\\
039 & 403,350 & 7,725,056\\
040 & 110,105 & 0\\
041 & 1,265,256 & 64,937,016\\
042 & 358,717 & 17,044,896\\
043 & 233,089 & 2,561,516\\
044 & 59,761 & 1,619,936\\
045 & 26,164 & 0\\
046 & 91,730 & 3,669,568\\
047 & 163,084 & 3,174,566\\
048 & 146,613 & 949,296\\
049 & 85,864 & 2,340,210\\
050 & 62,858 & 1,655,236\\
051 & 253,856 & 9,045,396\\
052 & 38,475 & 324,096\\
053 & 448,513 & 13,267,504\\
054 & 287,129 & 5,823,028\\
055 & 30,332 & 173,568\\
056 & 134,248 & 3,953,408\\
057 & 13,327 & 70,208\\
058 & 1,004,819 & 32,104,448\\
059 & 150,782 & 4,325,376\\
060 & 174,268 & 3,726,592\\
061 & 18,205 & 117,248\\
062 & 604,334 & 130,560\\
063 & 222,074 & 99,840\\
064 & 12,234 & 102,368\\
065 & 12,008 & 55,584\\
066 & 110,920 & 2,537,770\\
067 & 12,783,794 & 86,499,840\\
068 & 172,293 & 4,575,888\\
069 & 9,942 & 282,872\\
070 & 26,164 & 420,864\\
071 & 7,688 & 219,958\\
072 & 4,775,553 & 175,833,776\\
073 & 2,308,531 & 21,542,272\\
074 & 340,227 & 8,335,360\\
075 & 486,087 & 7,475,894\\
076 & 1,344,727 & 74,277,888\\
077 & 123,732 & 172,032\\
078 & 552,292 & 15,998,784\\
079 & 5,985,137 & 91,607,040\\
080 & 100,325 & 3,047,424\\
081 & 30,943 & 0\\
082 & 248,572 & 4,981,462\\
083 & 247,424 & 5,232,384\\
084 & 392,222 & 9,925,632\\
085 & 1,459,916 & 14,813,080\\
086 & 61,537 & 1,503,424\\
087 & 892,425 & 25,367,040\\
088 & 105,579 & 2,279,342\\
089 & 147,432 & 954,880\\
090 & 361,216 & 11,053,152\\
091 & 2,015,449 & 111,883,264\\
092 & 913,798 & 10,318,208\\
093 & 15,876 & 32,768\\
094 & 20,125 & 387,360\\
095 & 420,852 & 0\\
096 & 26,861 & 661,824\\
097 & 2,115,185 & 30,457,856\\
098 & 58,309 & 1,800,648\\*
099 & 338,418 & 8,744,376\\*
100 & 57,009 & 2,055,812\\
\end{longtable}
\endgroup

% 周期为总体任务执行时间；字节为额外 COPY。q1/q2 单核使用题设约定的原始整图响应。
\begingroup\normalsize
\setlength{\tabcolsep}{3pt}
\setlength{\LTleft}{\fill}\setlength{\LTright}{\fill}
\setlength{\LTcapwidth}{160mm}
\begin{longtable}{@{}p{15mm}>{\raggedleft\arraybackslash}p{28mm}>{\raggedleft\arraybackslash}p{39mm}>{\raggedleft\arraybackslash}p{28mm}>{\raggedleft\arraybackslash}p{39mm}@{}}
\caption{问题二逐例结果：1、2 核}\label{tab:appendix_q2_n12}\\
\toprule
算例 & \shortstack{1 核\\周期} & \shortstack{1 核额外\\COPY/B} & \shortstack{2 核\\周期} & \shortstack{2 核额外\\COPY/B}\\
\midrule
\endfirsthead
\caption[]{问题二逐例结果：1、2 核（续）}\\
\toprule
算例 & \shortstack{1 核\\周期} & \shortstack{1 核额外\\COPY/B} & \shortstack{2 核\\周期} & \shortstack{2 核额外\\COPY/B}\\
\midrule
\endhead
\bottomrule
\endfoot
001 & 233,110 & 0 & 116,868 & 1,152\\
002 & 261,945 & 0 & 127,390 & 1,181,184\\
003 & 602,212 & 0 & 304,083 & 2,966,846\\
004 & 392,357 & 0 & 199,074 & 110,592\\
005 & 95,618 & 0 & 49,416 & 944,678\\
006 & 82,837 & 16,400 & 37,602 & 128,396\\
007 & 1,012,562 & 0 & 490,814 & 27,648\\
008 & 487,605 & 0 & 136,146 & 3,538,944\\
009 & 153,261 & 0 & 83,434 & 44,262\\
010 & 86,969 & 0 & 37,534 & 241,840\\
011 & 184,697 & 0 & 92,646 & 180,224\\
012 & 59,088 & 0 & 27,188 & 49,888\\
013 & 263,949 & 0 & 132,198 & 8,192\\
014 & 17,698,626 & 71,010,720 & 8,829,656 & 72,145,816\\
015 & 177,362 & 180,256 & 78,032 & 109,824\\
016 & 7,715,523 & 0 & 4,026,536 & 404,188\\
017 & 162,810 & 22,528 & 74,648 & 379,288\\
018 & 1,004,367 & 1,970,176 & 497,540 & 3,483,648\\
019 & 81,144 & 0 & 36,508 & 263,716\\
020 & 532,746 & 0 & 266,580 & 0\\
021 & 6,567,937 & 4,501,504 & 3,163,679 & 8,820,736\\
022 & 56,780 & 0 & 27,917 & 150,276\\
023 & 122,614 & 0 & 49,766 & 679,872\\
024 & 2,555,343 & 0 & 1,334,756 & 396,844\\
025 & 4,863,026 & 0 & 2,351,860 & 27,648\\
026 & 99,203 & 0 & 44,447 & 264,116\\
027 & 2,363,952 & 1,400,832 & 1,117,457 & 2,162,688\\
028 & 43,836,040 & 277,412,256 & 20,338,661 & 47,163,648\\
029 & 37,555 & 0 & 18,871 & 120,968\\
030 & 2,808,034 & 5,519,360 & 1,401,638 & 10,001,408\\
031 & 1,414,369 & 3,491,328 & 709,457 & 5,109,760\\
032 & 41,443 & 0 & 17,286 & 6,216\\
033 & 1,774,196 & 4,592,640 & 891,424 & 6,516,768\\
034 & 842,996 & 1,079,832 & 391,515 & 2,095,616\\
035 & 169,571 & 0 & 92,022 & 189,696\\
036 & 931,510 & 0 & 466,068 & 1,152\\
037 & 210,001 & 0 & 105,374 & 204,800\\
038 & 1,364,555 & 5,627,904 & 677,436 & 6,262,788\\
039 & 1,435,262 & 7,929,856 & 704,087 & 9,191,424\\
040 & 253,454 & 0 & 130,452 & 1,402,390\\
041 & 6,003,957 & 30,898,176 & 2,993,408 & 26,663,040\\
042 & 1,686,740 & 6,629,376 & 839,320 & 6,838,560\\
043 & 677,656 & 49,160 & 326,409 & 539,648\\
044 & 154,407 & 4,438,112 & 56,402 & 1,062,496\\
045 & 125,614 & 0 & 63,059 & 0\\
046 & 276,455 & 3,176,288 & 106,876 & 930,400\\
047 & 326,508 & 0 & 165,302 & 1,412,774\\
048 & 222,220 & 0 & 135,424 & 726,470\\
049 & 196,276 & 0 & 108,567 & 744,682\\
050 & 149,690 & 0 & 79,272 & 457,678\\
051 & 607,628 & 0 & 318,596 & 393,884\\
052 & 180,301 & 0 & 80,154 & 1,119,386\\
053 & 1,472,686 & 5,405,536 & 707,951 & 4,566,686\\
054 & 815,051 & 0 & 440,915 & 4,374,794\\
055 & 118,835 & 69,632 & 55,868 & 78,348\\
056 & 277,285 & 0 & 163,006 & 721,074\\
057 & 59,726 & 94,216 & 29,387 & 20,256\\
058 & 4,539,717 & 28,696,576 & 2,252,009 & 29,773,824\\
059 & 829,473 & 1,691,648 & 369,318 & 360,448\\
060 & 888,789 & 2,068,520 & 436,355 & 2,357,504\\
061 & 77,191 & 0 & 38,453 & 34,560\\
062 & 2,854,809 & 0 & 1,428,964 & 75,264\\
063 & 1,070,397 & 0 & 536,109 & 75,264\\
064 & 19,116 & 0 & 12,171 & 40,038\\
065 & 53,068 & 0 & 22,719 & 204,496\\
066 & 359,070 & 0 & 198,351 & 1,795,116\\
067 & 63,834,834 & 264,449,408 & 30,330,574 & 33,761,664\\
068 & 362,095 & 0 & 191,696 & 2,000,138\\
069 & 25,144 & 0 & 13,821 & 118,438\\
070 & 117,300 & 0 & 51,604 & 168,468\\
071 & 18,919 & 0 & 9,537 & 148,644\\
072 & 23,897,828 & 108,613,680 & 11,835,911 & 112,186,968\\
073 & 10,081,626 & 141,570,432 & 4,145,708 & 4,983,168\\
074 & 1,682,239 & 6,070,280 & 823,784 & 7,430,180\\
075 & 1,441,809 & 0 & 706,423 & 3,238,406\\
076 & 6,185,113 & 28,280,832 & 3,084,437 & 26,401,792\\
077 & 495,320 & 0 & 251,045 & 231,272\\
078 & 2,296,143 & 12,644,568 & 1,049,840 & 2,678,592\\
079 & 31,052,251 & 32,448,512 & 15,424,280 & 31,780,864\\
080 & 292,138 & 548,864 & 145,475 & 1,667,072\\
081 & 151,795 & 0 & 76,026 & 1,392,980\\
082 & 721,041 & 0 & 350,242 & 1,557,030\\
083 & 1,057,993 & 13,934,688 & 411,067 & 2,396,768\\
084 & 2,507,412 & 0 & 686,336 & 9,925,632\\
085 & 3,187,899 & 0 & 1,573,135 & 7,725,538\\
086 & 118,403 & 0 & 60,928 & 1,243,142\\
087 & 3,974,071 & 7,902,720 & 1,976,432 & 9,768,960\\
088 & 243,623 & 0 & 134,488 & 1,382,482\\
089 & 693,092 & 0 & 328,935 & 898,368\\
090 & 1,178,363 & 21,233,592 & 463,105 & 3,174,576\\
091 & 9,186,769 & 53,947,008 & 4,527,122 & 51,685,016\\
092 & 4,815,554 & 66,000,096 & 2,007,977 & 6,087,264\\
093 & 74,205 & 0 & 37,827 & 8,192\\
094 & 94,892 & 0 & 38,892 & 115,488\\
095 & 2,084,853 & 0 & 573,690 & 16,588,800\\
096 & 84,194 & 0 & 37,660 & 208,776\\
097 & 11,491,509 & 10,584,064 & 5,394,735 & 58,413,056\\
098 & 177,430 & 0 & 83,291 & 586,952\\*
099 & 1,622,999 & 5,484,544 & 809,724 & 7,327,872\\*
100 & 151,223 & 0 & 63,475 & 223,792\\
\end{longtable}
\endgroup

% 周期为总体任务执行时间；字节为额外 COPY。q1/q2 单核使用题设约定的原始整图响应。
\begingroup\normalsize
\setlength{\tabcolsep}{3pt}
\setlength{\LTleft}{\fill}\setlength{\LTright}{\fill}
\setlength{\LTcapwidth}{160mm}
\begin{longtable}{@{}p{15mm}>{\raggedleft\arraybackslash}p{28mm}>{\raggedleft\arraybackslash}p{39mm}>{\raggedleft\arraybackslash}p{28mm}>{\raggedleft\arraybackslash}p{39mm}@{}}
\caption{问题二逐例结果：3、4 核}\label{tab:appendix_q2_n34}\\
\toprule
算例 & \shortstack{3 核\\周期} & \shortstack{3 核额外\\COPY/B} & \shortstack{4 核\\周期} & \shortstack{4 核额外\\COPY/B}\\
\midrule
\endfirsthead
\caption[]{问题二逐例结果：3、4 核（续）}\\
\toprule
算例 & \shortstack{3 核\\周期} & \shortstack{3 核额外\\COPY/B} & \shortstack{4 核\\周期} & \shortstack{4 核额外\\COPY/B}\\
\midrule
\endhead
\bottomrule
\endfoot
001 & 78,545 & 2,304 & 58,984 & 3,456\\
002 & 84,144 & 1,039,872 & 65,286 & 1,168,896\\
003 & 228,313 & 4,483,988 & 215,373 & 5,490,634\\
004 & 130,725 & 1,686,528 & 102,206 & 737,280\\
005 & 44,387 & 1,352,268 & 40,448 & 1,592,914\\
006 & 25,659 & 145,808 & 20,640 & 148,372\\
007 & 326,994 & 55,296 & 246,058 & 82,944\\
008 & 102,151 & 774,144 & 80,014 & 294,912\\
009 & 74,506 & 660,488 & 64,326 & 903,228\\
010 & 25,290 & 283,684 & 20,361 & 308,680\\
011 & 63,119 & 1,703,936 & 46,880 & 540,672\\
012 & 18,806 & 182,000 & 14,517 & 206,140\\
013 & 88,351 & 16,384 & 66,982 & 24,576\\
014 & 5,904,875 & 75,233,800 & 4,494,357 & 86,669,544\\
015 & 57,667 & 869,388 & 43,367 & 317,952\\
016 & 2,876,680 & 12,192 & 2,241,788 & 14,628\\
017 & 52,488 & 382,696 & 40,703 & 492,236\\
018 & 336,888 & 3,838,976 & 252,609 & 5,701,376\\
019 & 22,970 & 261,668 & 19,110 & 412,992\\
020 & 177,950 & 0 & 133,704 & 0\\
021 & 2,095,770 & 5,382,144 & 1,550,324 & 4,608,000\\
022 & 18,320 & 40,736 & 14,482 & 38,412\\
023 & 32,247 & 700,818 & 30,377 & 402,148\\
024 & 953,368 & 4,032 & 743,601 & 4,836\\
025 & 1,568,824 & 55,296 & 1,177,424 & 82,944\\
026 & 32,969 & 210,116 & 25,934 & 554,080\\
027 & 765,403 & 11,333,632 & 559,196 & 3,133,440\\
028 & 13,535,451 & 37,262,592 & 10,105,660 & 58,525,056\\
029 & 12,743 & 0 & 11,028 & 0\\
030 & 936,531 & 6,208,512 & 699,628 & 6,326,272\\
031 & 468,502 & 7,098,112 & 368,570 & 7,291,652\\
032 & 11,657 & 12,288 & 9,787 & 12,288\\
033 & 592,692 & 6,934,080 & 444,141 & 8,334,436\\
034 & 262,877 & 2,415,616 & 196,495 & 2,861,568\\
035 & 81,904 & 1,751,176 & 73,468 & 873,206\\
036 & 311,345 & 2,304 & 233,584 & 3,456\\
037 & 71,271 & 1,390,592 & 53,320 & 614,400\\
038 & 451,915 & 5,531,140 & 340,110 & 7,149,056\\
039 & 466,665 & 11,837,440 & 353,105 & 12,718,080\\
040 & 101,669 & 1,605,656 & 82,154 & 1,453,066\\
041 & 1,993,803 & 34,537,500 & 1,493,500 & 36,428,800\\
042 & 557,972 & 7,477,824 & 421,578 & 11,356,748\\
043 & 246,592 & 674,748 & 225,545 & 868,500\\
044 & 43,186 & 1,027,264 & 43,186 & 1,027,264\\
045 & 42,228 & 20,480 & 31,871 & 0\\
046 & 82,603 & 2,382,016 & 78,720 & 2,791,200\\
047 & 126,353 & 2,029,640 & 108,703 & 2,350,384\\
048 & 104,892 & 996,810 & 104,892 & 996,810\\
049 & 86,812 & 879,504 & 73,759 & 1,974,234\\
050 & 59,076 & 1,170,688 & 53,567 & 822,838\\
051 & 227,412 & 952 & 178,025 & 1,140\\
052 & 57,247 & 728,914 & 44,538 & 731,480\\
053 & 492,593 & 7,132,556 & 402,409 & 9,421,454\\
054 & 331,824 & 6,345,140 & 282,842 & 7,253,370\\
055 & 38,319 & 555,816 & 31,809 & 540,160\\
056 & 133,607 & 1,819,948 & 123,271 & 2,357,378\\
057 & 19,648 & 392,978 & 15,405 & 168,288\\
058 & 1,498,789 & 30,691,328 & 1,133,379 & 31,375,360\\
059 & 250,622 & 1,363,968 & 185,216 & 1,081,344\\
060 & 271,919 & 2,816,256 & 205,168 & 3,186,816\\
061 & 27,905 & 33,792 & 20,940 & 91,136\\
062 & 960,383 & 101,376 & 754,470 & 5,114,880\\
063 & 358,847 & 124,416 & 274,542 & 64,512\\
064 & 10,655 & 59,852 & 9,777 & 94,446\\
065 & 14,688 & 70,720 & 13,069 & 204,636\\
066 & 144,282 & 1,941,498 & 124,938 & 2,489,548\\
067 & 20,138,559 & 37,729,024 & 15,086,030 & 33,770,624\\
068 & 150,239 & 3,168,382 & 133,223 & 2,730,366\\
069 & 10,068 & 253,194 & 10,068 & 253,194\\
070 & 39,532 & 336,944 & 30,029 & 346,636\\
071 & 8,094 & 190,666 & 7,183 & 215,148\\
072 & 7,896,352 & 112,434,990 & 5,961,028 & 121,780,956\\
073 & 2,806,421 & 7,274,752 & 2,693,970 & 8,199,168\\
074 & 545,017 & 8,351,772 & 413,032 & 6,656,004\\
075 & 581,132 & 11,875,084 & 464,939 & 10,942,354\\
076 & 2,059,050 & 30,170,112 & 1,540,465 & 25,039,844\\
077 & 177,317 & 144,708 & 129,294 & 231,108\\
078 & 711,668 & 5,756,160 & 534,951 & 7,071,552\\
079 & 10,173,157 & 31,387,648 & 7,618,976 & 30,461,952\\
080 & 100,325 & 2,158,592 & 78,504 & 1,843,200\\
081 & 51,943 & 0 & 39,588 & 0\\
082 & 265,681 & 2,148,552 & 225,543 & 3,007,506\\
083 & 291,262 & 6,096,064 & 219,143 & 4,024,096\\
084 & 458,172 & 9,179,136 & 354,575 & 580,608\\
085 & 1,193,752 & 10,444,556 & 996,489 & 29,437,680\\
086 & 53,981 & 1,748,364 & 47,629 & 1,766,354\\
087 & 1,329,369 & 10,941,338 & 1,037,126 & 12,121,998\\
088 & 102,645 & 1,853,900 & 86,930 & 2,009,106\\
089 & 203,941 & 1,806,592 & 151,364 & 716,160\\
090 & 330,909 & 5,011,200 & 295,510 & 7,748,832\\
091 & 3,039,758 & 60,819,456 & 2,292,842 & 60,119,080\\
092 & 1,336,558 & 23,663,808 & 999,812 & 17,053,472\\
093 & 25,103 & 16,384 & 19,779 & 24,576\\
094 & 27,085 & 257,024 & 21,777 & 318,816\\
095 & 426,487 & 2,654,208 & 333,285 & 589,824\\
096 & 25,771 & 220,640 & 20,369 & 585,424\\
097 & 3,580,858 & 14,168,064 & 2,610,789 & 14,884,864\\
098 & 57,671 & 707,856 & 45,000 & 538,828\\*
099 & 541,752 & 7,857,924 & 406,659 & 9,291,392\\*
100 & 46,049 & 437,264 & 37,807 & 901,972\\
\end{longtable}
\endgroup

% 周期为总体任务执行时间；字节为额外 COPY。q1/q2 单核使用题设约定的原始整图响应。
\begingroup\normalsize
\setlength{\tabcolsep}{3pt}
\setlength{\LTleft}{\fill}\setlength{\LTright}{\fill}
\setlength{\LTcapwidth}{160mm}
\begin{longtable}{@{}p{15mm}>{\raggedleft\arraybackslash}p{28mm}>{\raggedleft\arraybackslash}p{39mm}@{}}
\caption{问题二逐例结果：5 核}\label{tab:appendix_q2_n5}\\
\toprule
算例 & \shortstack{5 核\\周期} & \shortstack{5 核额外\\COPY/B}\\
\midrule
\endfirsthead
\caption[]{问题二逐例结果：5 核（续）}\\
\toprule
算例 & \shortstack{5 核\\周期} & \shortstack{5 核额外\\COPY/B}\\
\midrule
\endhead
\bottomrule
\endfoot
001 & 47,502 & 4,608\\
002 & 54,463 & 1,125,888\\
003 & 204,960 & 6,224,082\\
004 & 82,819 & 663,552\\
005 & 39,482 & 1,565,906\\
006 & 19,090 & 115,624\\
007 & 197,052 & 110,592\\
008 & 65,706 & 294,912\\
009 & 60,567 & 1,061,928\\
010 & 20,127 & 331,080\\
011 & 39,329 & 1,337,344\\
012 & 13,786 & 50,112\\
013 & 54,363 & 32,768\\
014 & 3,637,679 & 92,883,656\\
015 & 40,072 & 432,640\\
016 & 2,088,400 & 515,880\\
017 & 33,136 & 593,196\\
018 & 202,146 & 6,639,620\\
019 & 16,248 & 101,472\\
020 & 107,187 & 20,480\\
021 & 1,245,796 & 7,518,208\\
022 & 11,604 & 6,916\\
023 & 24,541 & 391,900\\
024 & 695,220 & 499,556\\
025 & 943,120 & 110,592\\
026 & 25,078 & 437,916\\
027 & 459,633 & 18,331,648\\
028 & 8,079,303 & 56,498,688\\
029 & 9,383 & 0\\
030 & 560,497 & 7,489,536\\
031 & 308,178 & 7,907,738\\
032 & 8,671 & 71,752\\
033 & 351,905 & 8,951,424\\
034 & 157,197 & 3,510,272\\
035 & 67,441 & 1,000,940\\
036 & 187,182 & 4,608\\
037 & 43,013 & 901,120\\
038 & 272,867 & 7,391,236\\
039 & 308,250 & 9,654,272\\
040 & 74,356 & 1,458,238\\
041 & 1,202,403 & 25,396,360\\
042 & 348,719 & 8,089,152\\
043 & 203,748 & 5,983,346\\
044 & 43,186 & 1,027,264\\
045 & 25,770 & 40,960\\
046 & 69,268 & 2,469,536\\
047 & 105,952 & 2,927,124\\
048 & 99,368 & 1,264,536\\
049 & 66,733 & 1,937,700\\
050 & 48,084 & 1,283,572\\
051 & 169,451 & 493,392\\
052 & 37,466 & 324,096\\
053 & 382,931 & 11,974,532\\
054 & 243,665 & 7,991,196\\
055 & 26,186 & 395,272\\
056 & 108,501 & 2,892,630\\
057 & 12,956 & 282,176\\
058 & 955,671 & 15,409,152\\
059 & 149,793 & 5,701,632\\
060 & 173,991 & 3,576,064\\
061 & 17,011 & 104,960\\
062 & 587,941 & 89,088\\
063 & 221,161 & 70,656\\
064 & 9,597 & 103,282\\
065 & 11,968 & 62,792\\
066 & 106,287 & 2,557,124\\
067 & 12,255,440 & 26,076,672\\
068 & 132,986 & 3,232,496\\
069 & 8,320 & 312,914\\
070 & 26,164 & 420,864\\
071 & 7,183 & 215,148\\
072 & 4,781,658 & 128,357,638\\
073 & 2,130,376 & 14,611,840\\
074 & 328,436 & 10,383,396\\
075 & 386,155 & 12,416,696\\
076 & 1,237,970 & 32,310,272\\
077 & 120,675 & 1,602,662\\
078 & 471,520 & 10,292,736\\
079 & 6,007,992 & 33,882,112\\
080 & 71,369 & 2,273,280\\
081 & 30,943 & 0\\
082 & 193,581 & 3,895,832\\
083 & 199,734 & 5,736,832\\
084 & 285,686 & 322,560\\
085 & 897,331 & 37,614,932\\
086 & 46,397 & 1,815,608\\
087 & 864,851 & 12,687,110\\
088 & 77,808 & 2,034,172\\
089 & 121,742 & 995,840\\
090 & 272,200 & 8,762,496\\
091 & 1,850,544 & 64,694,276\\
092 & 796,666 & 15,358,080\\
093 & 15,876 & 32,768\\
094 & 18,767 & 332,064\\
095 & 269,789 & 147,456\\
096 & 18,920 & 435,404\\
097 & 2,101,059 & 17,844,224\\
098 & 39,049 & 716,040\\*
099 & 328,434 & 10,162,180\\*
100 & 34,446 & 795,816\\
\end{longtable}
\endgroup

% 周期为总体任务执行时间；字节为额外 COPY。q1/q2 单核使用题设约定的原始整图响应。
\begingroup\normalsize
\setlength{\tabcolsep}{3pt}
\setlength{\LTleft}{\fill}\setlength{\LTright}{\fill}
\setlength{\LTcapwidth}{160mm}
\begin{longtable}{@{}p{15mm}>{\raggedleft\arraybackslash}p{28mm}>{\raggedleft\arraybackslash}p{39mm}>{\raggedleft\arraybackslash}p{28mm}>{\raggedleft\arraybackslash}p{39mm}@{}}
\caption{问题三逐例结果：1、2 核}\label{tab:appendix_q3_n12}\\
\toprule
算例 & \shortstack{1 核\\周期} & \shortstack{1 核额外\\COPY/B} & \shortstack{2 核\\周期} & \shortstack{2 核额外\\COPY/B}\\
\midrule
\endfirsthead
\caption[]{问题三逐例结果：1、2 核（续）}\\
\toprule
算例 & \shortstack{1 核\\周期} & \shortstack{1 核额外\\COPY/B} & \shortstack{2 核\\周期} & \shortstack{2 核额外\\COPY/B}\\
\midrule
\endhead
\bottomrule
\endfoot
001 & 233,110 & 0 & 116,868 & 1,152\\
002 & 254,968 & 3,551,232 & 126,611 & 1,368,576\\
003 & 575,978 & 0 & 304,077 & 2,966,846\\
004 & 392,357 & 0 & 199,074 & 110,592\\
005 & 86,533 & 0 & 49,708 & 946,886\\
006 & 73,271 & 44,076 & 37,602 & 128,396\\
007 & 979,262 & 0 & 490,814 & 27,648\\
008 & 269,325 & 3,760,128 & 136,146 & 3,538,944\\
009 & 142,603 & 330,948 & 83,434 & 44,262\\
010 & 74,628 & 0 & 36,535 & 178,024\\
011 & 184,697 & 0 & 92,646 & 180,224\\
012 & 53,481 & 0 & 27,065 & 97,006\\
013 & 263,949 & 0 & 132,198 & 8,192\\
014 & 17,554,044 & 52,618,752 & 8,759,990 & 72,145,816\\
015 & 166,732 & 0 & 85,242 & 760,240\\
016 & 7,715,523 & 0 & 4,025,704 & 404,188\\
017 & 144,939 & 409,620 & 72,519 & 87,552\\
018 & 1,003,110 & 1,970,176 & 494,389 & 3,317,504\\
019 & 64,662 & 124,416 & 36,497 & 282,660\\
020 & 532,746 & 0 & 266,580 & 0\\
021 & 6,509,573 & 4,552,704 & 3,176,969 & 3,579,904\\
022 & 54,592 & 403,716 & 27,917 & 150,276\\
023 & 98,689 & 201,752 & 49,728 & 679,860\\
024 & 2,555,343 & 0 & 1,333,924 & 396,844\\
025 & 4,703,066 & 0 & 2,351,860 & 27,648\\
026 & 81,009 & 0 & 44,445 & 264,116\\
027 & 2,289,407 & 10,100,736 & 1,117,560 & 2,162,688\\
028 & 40,765,509 & 31,733,760 & 20,238,719 & 94,315,920\\
029 & 37,555 & 0 & 18,871 & 149,664\\
030 & 2,807,891 & 5,519,360 & 1,400,031 & 2,354,176\\
031 & 1,403,039 & 3,281,920 & 707,774 & 5,152,768\\
032 & 34,065 & 0 & 17,286 & 6,216\\
033 & 1,769,212 & 4,737,024 & 883,778 & 5,077,536\\
034 & 823,980 & 427,008 & 391,580 & 2,076,168\\
035 & 165,528 & 0 & 92,022 & 189,696\\
036 & 931,510 & 0 & 466,068 & 1,152\\
037 & 210,001 & 0 & 105,374 & 204,800\\
038 & 1,352,555 & 3,683,328 & 674,382 & 6,790,656\\
039 & 1,424,699 & 8,220,672 & 698,189 & 9,957,376\\
040 & 244,619 & 244,408 & 130,452 & 1,402,390\\
041 & 5,984,182 & 22,643,712 & 2,984,306 & 31,470,608\\
042 & 1,651,239 & 6,491,136 & 831,867 & 7,068,960\\
043 & 656,850 & 2,483,184 & 328,008 & 506,880\\
044 & 75,173 & 0 & 43,882 & 930,400\\
045 & 125,614 & 0 & 63,059 & 0\\
046 & 211,358 & 0 & 106,042 & 930,400\\
047 & 303,591 & 0 & 165,489 & 1,410,854\\
048 & 201,196 & 0 & 135,478 & 757,062\\
049 & 181,926 & 278,036 & 108,567 & 744,682\\
050 & 143,136 & 665,104 & 78,745 & 485,524\\
051 & 607,628 & 0 & 318,295 & 393,884\\
052 & 169,942 & 91,648 & 80,005 & 1,119,386\\
053 & 1,367,418 & 65,098 & 707,762 & 4,566,686\\
054 & 790,910 & 939,546 & 441,987 & 4,365,002\\
055 & 102,631 & 1,429,632 & 55,868 & 78,348\\
056 & 259,765 & 9,218 & 163,006 & 721,074\\
057 & 58,573 & 73,728 & 29,387 & 20,256\\
058 & 4,525,946 & 28,942,336 & 2,247,911 & 33,566,720\\
059 & 738,041 & 0 & 369,318 & 360,448\\
060 & 860,185 & 1,757,312 & 436,329 & 2,357,504\\
061 & 77,191 & 0 & 38,453 & 34,560\\
062 & 2,854,809 & 0 & 1,428,964 & 75,264\\
063 & 1,070,397 & 0 & 536,080 & 75,264\\
064 & 16,886 & 0 & 12,556 & 46,214\\
065 & 45,516 & 0 & 22,028 & 143,036\\
066 & 326,438 & 6,144 & 193,440 & 1,350,136\\
067 & 61,063,003 & 34,926,592 & 30,117,593 & 50,086,272\\
068 & 339,156 & 323,896 & 191,254 & 2,000,138\\
069 & 21,899 & 0 & 13,816 & 118,438\\
070 & 109,766 & 0 & 51,503 & 168,468\\
071 & 16,084 & 0 & 9,783 & 244,966\\
072 & 23,669,887 & 89,090,992 & 11,807,422 & 104,426,272\\
073 & 8,221,623 & 188,416 & 4,184,375 & 8,734,592\\
074 & 1,644,373 & 6,479,872 & 813,107 & 5,988,360\\
075 & 1,358,701 & 0 & 706,423 & 3,238,406\\
076 & 6,153,011 & 15,784,960 & 3,082,459 & 26,401,792\\
077 & 488,026 & 0 & 248,734 & 57,344\\
078 & 2,067,989 & 0 & 1,024,634 & 2,148,480\\
079 & 30,856,884 & 32,772,096 & 15,388,269 & 31,875,072\\
080 & 290,734 & 1,040,384 & 145,338 & 1,785,856\\
081 & 151,795 & 0 & 76,026 & 1,392,980\\
082 & 678,162 & 538,506 & 348,070 & 1,623,014\\
083 & 841,782 & 131,072 & 410,764 & 5,288,544\\
084 & 1,371,156 & 9,676,800 & 686,336 & 9,925,632\\
085 & 3,097,342 & 3,983,714 & 1,573,135 & 7,725,538\\
086 & 108,508 & 72,390 & 61,165 & 1,253,446\\
087 & 3,948,565 & 8,317,770 & 1,979,975 & 9,768,960\\
088 & 231,497 & 391,168 & 131,063 & 1,343,442\\
089 & 593,264 & 0 & 295,856 & 238,720\\
090 & 843,811 & 0 & 451,406 & 3,341,952\\
091 & 9,112,016 & 36,733,952 & 4,527,598 & 49,288,192\\
092 & 4,010,076 & 5,476,352 & 2,009,860 & 6,070,880\\
093 & 74,205 & 0 & 37,827 & 8,192\\
094 & 86,556 & 0 & 38,992 & 115,200\\
095 & 1,144,413 & 16,809,984 & 573,690 & 16,588,800\\
096 & 73,473 & 22,584 & 36,346 & 196,484\\
097 & 10,934,396 & 58,613,760 & 5,398,228 & 58,441,728\\
\noalign{\penalty-10000}
098 & 166,777 & 372,804 & 83,520 & 550,092\\*
099 & 1,609,527 & 4,550,656 & 806,783 & 7,327,872\\*
100 & 122,765 & 278,128 & 64,209 & 624,972\\
\end{longtable}
\endgroup

% 周期为总体任务执行时间；字节为额外 COPY。q1/q2 单核使用题设约定的原始整图响应。
\begingroup\normalsize
\setlength{\tabcolsep}{3pt}
\setlength{\LTleft}{\fill}\setlength{\LTright}{\fill}
\setlength{\LTcapwidth}{160mm}
\begin{longtable}{@{}p{15mm}>{\raggedleft\arraybackslash}p{28mm}>{\raggedleft\arraybackslash}p{39mm}>{\raggedleft\arraybackslash}p{28mm}>{\raggedleft\arraybackslash}p{39mm}@{}}
\caption{问题三逐例结果：3、4 核}\label{tab:appendix_q3_n34}\\
\toprule
算例 & \shortstack{3 核\\周期} & \shortstack{3 核额外\\COPY/B} & \shortstack{4 核\\周期} & \shortstack{4 核额外\\COPY/B}\\
\midrule
\endfirsthead
\caption[]{问题三逐例结果：3、4 核（续）}\\
\toprule
算例 & \shortstack{3 核\\周期} & \shortstack{3 核额外\\COPY/B} & \shortstack{4 核\\周期} & \shortstack{4 核额外\\COPY/B}\\
\midrule
\endhead
\bottomrule
\endfoot
001 & 78,545 & 2,304 & 58,984 & 3,456\\
002 & 84,144 & 1,039,872 & 65,286 & 1,168,896\\
003 & 200,752 & 1,957,280 & 177,413 & 3,522,458\\
004 & 131,412 & 1,723,392 & 102,206 & 724,992\\
005 & 42,683 & 1,392,620 & 37,678 & 1,574,162\\
006 & 25,659 & 145,808 & 21,977 & 144,788\\
007 & 326,994 & 55,296 & 246,058 & 82,944\\
008 & 102,151 & 774,144 & 80,014 & 294,912\\
009 & 70,133 & 785,740 & 63,026 & 850,506\\
010 & 24,221 & 219,688 & 21,442 & 954,200\\
011 & 62,193 & 3,414,016 & 46,808 & 704,512\\
012 & 18,828 & 173,808 & 14,371 & 159,024\\
013 & 88,351 & 16,384 & 66,946 & 24,576\\
014 & 5,845,375 & 80,144,392 & 4,399,727 & 60,560,280\\
015 & 59,525 & 914,160 & 45,687 & 317,952\\
016 & 2,876,680 & 12,192 & 2,241,788 & 14,628\\
017 & 49,218 & 175,104 & 38,361 & 253,440\\
018 & 332,384 & 3,741,696 & 251,720 & 3,791,872\\
019 & 22,970 & 287,524 & 18,443 & 380,228\\
020 & 177,950 & 0 & 133,704 & 0\\
021 & 2,109,297 & 5,410,816 & 1,550,324 & 4,608,000\\
022 & 18,743 & 23,328 & 14,388 & 38,412\\
023 & 32,743 & 622,102 & 29,297 & 1,306,496\\
024 & 953,368 & 4,032 & 743,601 & 4,836\\
025 & 1,568,824 & 55,296 & 1,177,424 & 82,944\\
026 & 32,360 & 654,940 & 26,181 & 511,592\\
027 & 758,033 & 11,167,744 & 559,196 & 3,133,440\\
028 & 13,541,489 & 37,996,800 & 10,063,752 & 94,118,784\\
029 & 12,743 & 0 & 10,424 & 280,968\\
030 & 931,803 & 4,257,792 & 696,606 & 5,032,960\\
031 & 471,340 & 3,615,744 & 355,023 & 5,144,064\\
032 & 11,775 & 12,288 & 9,481 & 12,300\\
033 & 587,082 & 7,426,184 & 442,630 & 7,214,688\\
034 & 265,966 & 2,320,384 & 196,515 & 2,861,568\\
035 & 79,468 & 584,908 & 68,430 & 1,097,894\\
036 & 311,345 & 2,304 & 233,584 & 3,456\\
037 & 71,193 & 1,390,592 & 53,320 & 614,400\\
038 & 451,791 & 5,426,692 & 338,219 & 5,772,288\\
039 & 468,365 & 11,755,520 & 349,027 & 12,656,640\\
040 & 101,669 & 1,605,656 & 85,446 & 1,720,946\\
041 & 1,988,631 & 35,111,940 & 1,490,581 & 36,428,816\\
042 & 552,280 & 7,577,152 & 420,768 & 11,353,676\\
043 & 246,359 & 674,748 & 221,351 & 868,500\\
044 & 36,208 & 2,194,368 & 36,208 & 2,194,368\\
045 & 42,228 & 20,480 & 31,871 & 0\\
046 & 77,417 & 2,220,224 & 62,397 & 2,905,888\\
047 & 126,167 & 2,029,640 & 108,029 & 2,337,904\\
048 & 104,892 & 996,810 & 104,892 & 996,810\\
049 & 83,117 & 1,008,040 & 71,220 & 1,837,534\\
050 & 59,063 & 1,204,662 & 51,949 & 1,540,648\\
051 & 227,412 & 952 & 178,025 & 1,140\\
052 & 57,447 & 1,282,654 & 44,460 & 632,750\\
053 & 485,409 & 7,234,968 & 384,510 & 9,301,518\\
054 & 332,438 & 6,193,716 & 286,390 & 7,240,730\\
055 & 37,611 & 115,712 & 30,914 & 580,372\\
056 & 132,899 & 1,819,948 & 121,253 & 2,402,690\\
057 & 19,648 & 392,978 & 15,544 & 447,792\\
058 & 1,495,460 & 31,043,584 & 1,118,279 & 31,543,296\\
059 & 250,289 & 1,363,968 & 185,216 & 1,081,344\\
060 & 268,559 & 2,908,416 & 207,019 & 3,093,632\\
061 & 27,905 & 33,792 & 20,940 & 91,136\\
062 & 963,935 & 73,728 & 750,199 & 22,798,848\\
063 & 358,787 & 55,296 & 274,542 & 64,512\\
064 & 10,820 & 59,916 & 9,629 & 94,446\\
065 & 14,513 & 85,056 & 11,713 & 119,360\\
066 & 138,440 & 2,665,988 & 117,281 & 2,438,828\\
067 & 19,951,926 & 58,923,776 & 14,997,101 & 56,276,096\\
068 & 147,114 & 3,326,302 & 127,752 & 2,789,758\\
069 & 10,008 & 253,194 & 10,008 & 253,194\\
070 & 39,077 & 345,136 & 29,793 & 346,636\\
071 & 8,430 & 227,976 & 7,102 & 227,436\\
072 & 7,866,605 & 112,434,990 & 5,924,972 & 122,580,790\\
073 & 2,759,896 & 7,844,608 & 2,125,171 & 11,164,800\\
074 & 544,852 & 8,351,772 & 405,974 & 7,229,452\\
075 & 579,216 & 11,876,108 & 460,158 & 10,952,338\\
076 & 2,047,158 & 18,502,656 & 1,529,141 & 20,473,856\\
077 & 169,489 & 114,688 & 126,806 & 114,688\\
078 & 701,842 & 7,239,360 & 517,424 & 6,973,824\\
079 & 10,113,306 & 38,875,136 & 7,550,155 & 42,356,736\\
080 & 100,325 & 2,285,568 & 76,782 & 1,708,032\\
081 & 51,943 & 0 & 39,588 & 0\\
082 & 265,681 & 2,148,552 & 226,842 & 3,139,858\\
083 & 281,995 & 3,277,760 & 211,087 & 2,791,200\\
084 & 458,172 & 9,179,136 & 353,412 & 737,280\\
085 & 1,193,708 & 10,444,556 & 986,045 & 29,437,680\\
086 & 51,974 & 1,759,116 & 44,949 & 2,088,146\\
087 & 1,345,661 & 10,874,522 & 1,045,358 & 9,695,232\\
088 & 99,048 & 1,895,308 & 85,617 & 1,807,122\\
089 & 203,894 & 1,806,592 & 148,856 & 716,160\\
090 & 318,468 & 5,760,384 & 255,234 & 6,606,720\\
091 & 3,032,804 & 60,819,456 & 2,280,824 & 65,764,500\\
092 & 1,339,219 & 7,521,472 & 996,428 & 6,264,608\\
093 & 25,103 & 16,384 & 19,743 & 24,576\\
094 & 27,531 & 211,968 & 21,738 & 318,816\\
095 & 433,227 & 2,433,024 & 335,913 & 147,840\\
096 & 26,191 & 290,508 & 19,873 & 649,432\\
097 & 3,551,104 & 14,641,152 & 2,596,536 & 17,580,032\\
098 & 57,758 & 707,860 & 46,309 & 784,600\\*
099 & 540,097 & 8,877,696 & 403,923 & 9,291,392\\*
100 & 46,369 & 570,616 & 38,002 & 1,283,016\\
\end{longtable}
\endgroup

% 周期为总体任务执行时间；字节为额外 COPY。q1/q2 单核使用题设约定的原始整图响应。
\begingroup\normalsize
\setlength{\tabcolsep}{3pt}
\setlength{\LTleft}{\fill}\setlength{\LTright}{\fill}
\setlength{\LTcapwidth}{160mm}
\begin{longtable}{@{}p{15mm}>{\raggedleft\arraybackslash}p{28mm}>{\raggedleft\arraybackslash}p{39mm}@{}}
\caption{问题三逐例结果：5 核}\label{tab:appendix_q3_n5}\\
\toprule
算例 & \shortstack{5 核\\周期} & \shortstack{5 核额外\\COPY/B}\\
\midrule
\endfirsthead
\caption[]{问题三逐例结果：5 核（续）}\\
\toprule
算例 & \shortstack{5 核\\周期} & \shortstack{5 核额外\\COPY/B}\\
\midrule
\endhead
\bottomrule
\endfoot
001 & 47,502 & 4,608\\
002 & 53,957 & 1,230,336\\
003 & 142,684 & 3,900,208\\
004 & 82,819 & 651,264\\
005 & 37,678 & 1,574,162\\
006 & 18,052 & 210,348\\
007 & 197,052 & 110,592\\
008 & 65,706 & 294,912\\
009 & 54,001 & 988,306\\
010 & 17,155 & 292,560\\
011 & 39,393 & 1,447,936\\
012 & 12,615 & 251,716\\
013 & 54,117 & 32,768\\
014 & 3,535,676 & 89,325,084\\
015 & 38,817 & 432,640\\
016 & 2,086,586 & 515,880\\
017 & 31,159 & 336,384\\
018 & 200,960 & 3,120,000\\
019 & 16,415 & 98,400\\
020 & 107,187 & 20,480\\
021 & 1,247,880 & 9,588,736\\
022 & 11,474 & 13,824\\
023 & 24,350 & 398,044\\
024 & 693,470 & 499,552\\
025 & 943,120 & 110,592\\
026 & 23,047 & 453,240\\
027 & 452,503 & 17,725,440\\
028 & 8,041,278 & 34,406,400\\
029 & 9,383 & 0\\
030 & 558,512 & 6,318,080\\
031 & 284,116 & 5,666,304\\
032 & 8,396 & 59,520\\
033 & 350,553 & 8,951,424\\
034 & 157,798 & 3,510,272\\
035 & 65,383 & 718,326\\
036 & 187,182 & 4,608\\
037 & 43,013 & 901,120\\
038 & 272,000 & 7,279,108\\
039 & 279,449 & 14,704,640\\
040 & 75,955 & 1,592,310\\
041 & 1,194,456 & 40,669,188\\
042 & 339,457 & 10,187,328\\
043 & 206,007 & 6,002,246\\
044 & 36,208 & 2,194,368\\
045 & 25,770 & 40,960\\
046 & 49,119 & 3,527,040\\
047 & 102,266 & 2,927,124\\
048 & 98,766 & 1,264,536\\
049 & 67,575 & 2,011,460\\
050 & 47,546 & 1,520,298\\
051 & 167,637 & 493,392\\
052 & 37,287 & 950,052\\
053 & 361,811 & 12,007,080\\
054 & 242,741 & 7,988,828\\
055 & 23,160 & 173,568\\
056 & 107,262 & 2,832,584\\
057 & 12,956 & 216,640\\
058 & 911,574 & 42,000,384\\
059 & 150,246 & 5,713,920\\
060 & 173,733 & 3,576,064\\
061 & 16,981 & 104,960\\
062 & 587,941 & 82,944\\
063 & 221,161 & 67,584\\
064 & 9,227 & 95,122\\
065 & 10,267 & 74,160\\
066 & 104,620 & 2,622,404\\
067 & 12,061,334 & 71,656,960\\
068 & 121,498 & 3,022,320\\
069 & 7,886 & 320,114\\
070 & 24,137 & 435,200\\
071 & 5,619 & 272,274\\
072 & 4,759,988 & 131,848,768\\
073 & 1,738,206 & 20,287,488\\
074 & 330,165 & 12,501,052\\
075 & 369,043 & 12,400,568\\
076 & 1,224,141 & 23,614,464\\
077 & 115,617 & 172,032\\
078 & 433,782 & 12,005,376\\
079 & 6,062,064 & 33,882,112\\
080 & 63,001 & 3,616,768\\
081 & 30,943 & 0\\
082 & 184,003 & 3,820,568\\
083 & 177,330 & 5,802,368\\
084 & 285,686 & 322,560\\
085 & 862,197 & 37,590,868\\
086 & 44,818 & 2,088,146\\
087 & 829,663 & 12,487,680\\
088 & 74,720 & 2,089,020\\
089 & 120,170 & 954,880\\
090 & 196,502 & 13,899,264\\
091 & 1,828,683 & 43,944,448\\
092 & 773,666 & 6,433,152\\
093 & 15,595 & 36,864\\
094 & 18,035 & 332,064\\
095 & 269,789 & 147,456\\
096 & 18,630 & 442,444\\
097 & 2,100,021 & 18,178,048\\
098 & 39,669 & 830,752\\*
099 & 324,499 & 10,166,276\\*
100 & 32,716 & 777,764\\
\end{longtable}
\endgroup

% 由正式500组同方案完整评价记录生成；仅显示命中率保留六位小数。
\Needspace{8\baselineskip}
\begingroup\normalsize
\setlength{\tabcolsep}{1.5pt}
\setlength{\LTleft}{\fill}\setlength{\LTright}{\fill}
\setlength{\LTcapwidth}{160mm}
\begin{longtable}{@{}p{12mm}>{\raggedleft\arraybackslash}p{9mm}>{\raggedleft\arraybackslash}p{24mm}>{\raggedleft\arraybackslash}p{24mm}>{\raggedleft\arraybackslash}p{30mm}>{\raggedleft\arraybackslash}p{30mm}>{\raggedleft\arraybackslash}p{23mm}@{}}
\caption{问题三逐图逐核同方案有无 L2 的官方评价结果}\label{tab:appendix_q3_all_pairs}\\
\toprule
算例 & 核数 & \shortstack{无L2\\周期} & \shortstack{有L2\\周期} & \shortstack{无L2额外\\COPY/B} & \shortstack{有L2额外\\COPY/B} & \shortstack{字节\\命中率}\\
\midrule
\endfirsthead
\caption[]{问题三逐图逐核同方案有无L2的官方评价结果（续）}\\
\toprule
算例 & 核数 & \shortstack{无L2\\周期} & \shortstack{有L2\\周期} & \shortstack{无L2额外\\COPY/B} & \shortstack{有L2额外\\COPY/B} & \shortstack{字节\\命中率}\\
\midrule
\endhead
\bottomrule
\endfoot
001 & 1 & 233,110 & 233,110 & 0 & 0 & 0.000000\\
001 & 2 & 116,868 & 116,868 & 1,152 & 1,152 & 0.000000\\
001 & 3 & 78,545 & 78,545 & 2,304 & 2,304 & 0.000000\\
001 & 4 & 58,984 & 58,984 & 3,456 & 3,456 & 0.000000\\
001 & 5 & 47,502 & 47,502 & 4,608 & 4,608 & 0.000000\\
002 & 1 & 255,084 & 254,968 & 3,551,232 & 3,551,232 & 0.345979\\
002 & 2 & 126,699 & 126,611 & 1,368,576 & 1,368,576 & 0.298149\\
002 & 3 & 84,144 & 84,144 & 1,039,872 & 1,039,872 & 0.295267\\
002 & 4 & 65,286 & 65,286 & 1,168,896 & 1,168,896 & 0.253400\\
002 & 5 & 53,957 & 53,957 & 1,230,336 & 1,230,336 & 0.325402\\
003 & 1 & 575,978 & 575,978 & 0 & 0 & 0.000000\\
003 & 2 & 304,083 & 304,077 & 2,966,846 & 2,966,846 & 0.134344\\
003 & 3 & 200,759 & 200,752 & 1,957,280 & 1,957,280 & 0.308343\\
003 & 4 & 180,000 & 177,413 & 3,522,458 & 3,522,458 & 0.408295\\
003 & 5 & 144,728 & 142,684 & 3,900,208 & 3,900,208 & 0.500822\\
004 & 1 & 392,357 & 392,357 & 0 & 0 & 0.000000\\
004 & 2 & 199,074 & 199,074 & 110,592 & 110,592 & 0.000000\\
004 & 3 & 131,412 & 131,412 & 1,723,392 & 1,723,392 & 0.160714\\
004 & 4 & 102,206 & 102,206 & 724,992 & 724,992 & 0.012461\\
004 & 5 & 82,819 & 82,819 & 651,264 & 651,264 & 0.060790\\
005 & 1 & 86,533 & 86,533 & 0 & 0 & 0.000000\\
005 & 2 & 49,833 & 49,708 & 946,886 & 946,886 & 0.083254\\
005 & 3 & 43,953 & 42,683 & 1,392,620 & 1,392,620 & 0.300078\\
005 & 4 & 40,257 & 37,678 & 1,574,162 & 1,574,162 & 0.320578\\
005 & 5 & 40,257 & 37,678 & 1,574,162 & 1,574,162 & 0.320578\\
006 & 1 & 73,271 & 73,271 & 44,076 & 44,076 & 0.000000\\
006 & 2 & 37,602 & 37,602 & 128,396 & 128,396 & 0.191146\\
006 & 3 & 25,659 & 25,659 & 145,808 & 145,808 & 0.211580\\
006 & 4 & 21,979 & 21,977 & 144,788 & 144,788 & 0.197011\\
006 & 5 & 18,052 & 18,052 & 210,348 & 210,348 & 0.266843\\
007 & 1 & 979,262 & 979,262 & 0 & 0 & 0.000000\\
007 & 2 & 490,814 & 490,814 & 27,648 & 27,648 & 0.000000\\
007 & 3 & 326,994 & 326,994 & 55,296 & 55,296 & 0.000000\\
007 & 4 & 246,058 & 246,058 & 82,944 & 82,944 & 0.000000\\
007 & 5 & 197,052 & 197,052 & 110,592 & 110,592 & 0.000000\\
008 & 1 & 269,325 & 269,325 & 3,760,128 & 3,760,128 & 0.000000\\
008 & 2 & 136,146 & 136,146 & 3,538,944 & 3,538,944 & 0.000000\\
008 & 3 & 102,151 & 102,151 & 774,144 & 774,144 & 0.000000\\
008 & 4 & 80,014 & 80,014 & 294,912 & 294,912 & 0.000000\\
008 & 5 & 65,706 & 65,706 & 294,912 & 294,912 & 0.000000\\
009 & 1 & 142,869 & 142,603 & 330,948 & 330,948 & 0.157107\\
009 & 2 & 83,434 & 83,434 & 44,262 & 44,262 & 0.013833\\
009 & 3 & 70,134 & 70,133 & 785,740 & 785,740 & 0.149949\\
009 & 4 & 63,131 & 63,026 & 850,506 & 850,506 & 0.209327\\
009 & 5 & 55,040 & 54,001 & 988,306 & 988,306 & 0.223435\\
010 & 1 & 74,628 & 74,628 & 0 & 0 & 0.000000\\
010 & 2 & 36,535 & 36,535 & 178,024 & 178,024 & 0.208651\\
010 & 3 & 25,262 & 24,221 & 219,688 & 219,688 & 0.285099\\
010 & 4 & 26,178 & 21,442 & 954,200 & 954,200 & 0.334449\\
010 & 5 & 18,311 & 17,155 & 292,560 & 292,560 & 0.408809\\
011 & 1 & 184,697 & 184,697 & 0 & 0 & 0.000000\\
011 & 2 & 92,646 & 92,646 & 180,224 & 180,224 & 0.000000\\
011 & 3 & 67,891 & 62,193 & 3,414,016 & 3,414,016 & 0.271923\\
011 & 4 & 46,880 & 46,808 & 704,512 & 704,512 & 0.135714\\
011 & 5 & 39,906 & 39,393 & 1,447,936 & 1,447,936 & 0.288344\\
012 & 1 & 53,481 & 53,481 & 0 & 0 & 0.000000\\
012 & 2 & 27,168 & 27,065 & 97,006 & 97,006 & 0.140652\\
012 & 3 & 18,828 & 18,828 & 173,808 & 173,808 & 0.109356\\
012 & 4 & 14,377 & 14,371 & 159,024 & 159,024 & 0.141160\\
012 & 5 & 14,162 & 12,615 & 251,716 & 251,716 & 0.152586\\
013 & 1 & 263,949 & 263,949 & 0 & 0 & 0.000000\\
013 & 2 & 132,198 & 132,198 & 8,192 & 8,192 & 0.000000\\
013 & 3 & 88,351 & 88,351 & 16,384 & 16,384 & 0.000000\\
013 & 4 & 66,982 & 66,946 & 24,576 & 24,576 & 0.009000\\
013 & 5 & 54,363 & 54,117 & 32,768 & 32,768 & 0.011952\\
014 & 1 & 17,620,903 & 17,554,044 & 52,618,752 & 52,618,752 & 0.228455\\
014 & 2 & 8,829,656 & 8,759,990 & 72,145,816 & 72,145,816 & 0.301520\\
014 & 3 & 5,915,638 & 5,845,375 & 80,144,392 & 80,144,392 & 0.361812\\
014 & 4 & 4,401,505 & 4,399,727 & 60,560,280 & 60,560,280 & 0.068330\\
014 & 5 & 3,617,022 & 3,535,676 & 89,325,084 & 89,325,084 & 0.424704\\
015 & 1 & 166,732 & 166,732 & 0 & 0 & 0.000000\\
015 & 2 & 85,544 & 85,242 & 760,240 & 760,240 & 0.325333\\
015 & 3 & 61,214 & 59,525 & 914,160 & 914,160 & 0.333368\\
015 & 4 & 46,139 & 45,687 & 317,952 & 317,952 & 0.209921\\
015 & 5 & 40,072 & 38,817 & 432,640 & 432,640 & 0.154324\\
016 & 1 & 7,715,523 & 7,715,523 & 0 & 0 & 0.000000\\
016 & 2 & 4,026,536 & 4,025,704 & 404,188 & 404,188 & 0.496536\\
016 & 3 & 2,876,680 & 2,876,680 & 12,192 & 12,192 & 0.001523\\
016 & 4 & 2,241,788 & 2,241,788 & 14,628 & 14,628 & 0.003036\\
016 & 5 & 2,088,400 & 2,086,586 & 515,880 & 515,880 & 0.550045\\
017 & 1 & 145,165 & 144,939 & 409,620 & 409,620 & 0.192600\\
017 & 2 & 72,529 & 72,519 & 87,552 & 87,552 & 0.022201\\
017 & 3 & 49,218 & 49,218 & 175,104 & 175,104 & 0.246548\\
017 & 4 & 38,367 & 38,361 & 253,440 & 253,440 & 0.312631\\
017 & 5 & 31,159 & 31,159 & 336,384 & 336,384 & 0.047587\\
018 & 1 & 1,004,367 & 1,003,110 & 1,970,176 & 1,970,176 & 0.343750\\
018 & 2 & 495,014 & 494,389 & 3,317,504 & 3,317,504 & 0.443038\\
018 & 3 & 332,384 & 332,384 & 3,741,696 & 3,741,696 & 0.499302\\
018 & 4 & 252,966 & 251,720 & 3,791,872 & 3,791,872 & 0.653494\\
018 & 5 & 200,960 & 200,960 & 3,120,000 & 3,120,000 & 0.591076\\
019 & 1 & 64,662 & 64,662 & 124,416 & 124,416 & 0.081367\\
019 & 2 & 36,497 & 36,497 & 282,660 & 282,660 & 0.171645\\
019 & 3 & 22,970 & 22,970 & 287,524 & 287,524 & 0.226396\\
019 & 4 & 18,443 & 18,443 & 380,228 & 380,228 & 0.213263\\
019 & 5 & 16,415 & 16,415 & 98,400 & 98,400 & 0.166301\\
020 & 1 & 532,746 & 532,746 & 0 & 0 & 0.000000\\
020 & 2 & 266,580 & 266,580 & 0 & 0 & 0.000000\\
020 & 3 & 177,950 & 177,950 & 0 & 0 & 0.000000\\
020 & 4 & 133,704 & 133,704 & 0 & 0 & 0.000000\\
020 & 5 & 107,187 & 107,187 & 20,480 & 20,480 & 0.003436\\
021 & 1 & 6,517,945 & 6,509,573 & 4,552,704 & 4,552,704 & 0.486511\\
021 & 2 & 3,177,619 & 3,176,969 & 3,579,904 & 3,579,904 & 0.489303\\
021 & 3 & 2,111,226 & 2,109,297 & 5,410,816 & 5,410,816 & 0.558761\\
021 & 4 & 1,550,324 & 1,550,324 & 4,608,000 & 4,608,000 & 0.000000\\
021 & 5 & 1,248,347 & 1,247,880 & 9,588,736 & 9,588,736 & 0.601509\\
022 & 1 & 54,592 & 54,592 & 403,716 & 403,716 & 0.019528\\
022 & 2 & 27,917 & 27,917 & 150,276 & 150,276 & 0.023660\\
022 & 3 & 18,743 & 18,743 & 23,328 & 23,328 & 0.090133\\
022 & 4 & 14,388 & 14,388 & 38,412 & 38,412 & 0.109613\\
022 & 5 & 11,474 & 11,474 & 13,824 & 13,824 & 0.060801\\
023 & 1 & 98,689 & 98,689 & 201,752 & 201,752 & 0.003821\\
023 & 2 & 49,728 & 49,728 & 679,860 & 679,860 & 0.076508\\
023 & 3 & 33,014 & 32,743 & 622,102 & 622,102 & 0.148585\\
023 & 4 & 33,600 & 29,297 & 1,306,496 & 1,306,496 & 0.244014\\
023 & 5 & 24,649 & 24,350 & 398,044 & 398,044 & 0.192048\\
024 & 1 & 2,555,343 & 2,555,343 & 0 & 0 & 0.000000\\
024 & 2 & 1,334,756 & 1,333,924 & 396,844 & 396,844 & 0.498849\\
024 & 3 & 953,368 & 953,368 & 4,032 & 4,032 & 0.000506\\
024 & 4 & 743,601 & 743,601 & 4,836 & 4,836 & 0.001011\\
024 & 5 & 695,284 & 693,470 & 499,552 & 499,552 & 0.553720\\
025 & 1 & 4,703,066 & 4,703,066 & 0 & 0 & 0.000000\\
025 & 2 & 2,351,860 & 2,351,860 & 27,648 & 27,648 & 0.000000\\
025 & 3 & 1,568,824 & 1,568,824 & 55,296 & 55,296 & 0.000000\\
025 & 4 & 1,177,424 & 1,177,424 & 82,944 & 82,944 & 0.000000\\
025 & 5 & 943,120 & 943,120 & 110,592 & 110,592 & 0.000000\\
026 & 1 & 81,009 & 81,009 & 0 & 0 & 0.000000\\
026 & 2 & 44,447 & 44,445 & 264,116 & 264,116 & 0.346196\\
026 & 3 & 33,706 & 32,360 & 654,940 & 654,940 & 0.298733\\
026 & 4 & 26,970 & 26,181 & 511,592 & 511,592 & 0.388758\\
026 & 5 & 25,762 & 23,047 & 453,240 & 453,240 & 0.378393\\
027 & 1 & 2,289,407 & 2,289,407 & 10,100,736 & 10,100,736 & 0.000000\\
027 & 2 & 1,117,560 & 1,117,560 & 2,162,688 & 2,162,688 & 0.171779\\
027 & 3 & 761,877 & 758,033 & 11,167,744 & 11,167,744 & 0.294904\\
027 & 4 & 559,196 & 559,196 & 3,133,440 & 3,133,440 & 0.210700\\
027 & 5 & 455,659 & 452,503 & 17,725,440 & 17,725,440 & 0.285102\\
028 & 1 & 40,765,509 & 40,765,509 & 31,733,760 & 31,733,760 & 0.000000\\
028 & 2 & 20,243,537 & 20,238,719 & 94,315,920 & 94,315,920 & 0.038358\\
028 & 3 & 13,554,271 & 13,541,489 & 37,996,800 & 37,996,800 & 0.138807\\
028 & 4 & 10,087,972 & 10,063,752 & 94,118,784 & 94,118,784 & 0.107932\\
028 & 5 & 8,071,642 & 8,041,278 & 34,406,400 & 34,406,400 & 0.334632\\
029 & 1 & 37,555 & 37,555 & 0 & 0 & 0.000000\\
029 & 2 & 18,871 & 18,871 & 149,664 & 149,664 & 0.343296\\
029 & 3 & 12,743 & 12,743 & 0 & 0 & 0.000000\\
029 & 4 & 12,480 & 10,424 & 280,968 & 280,968 & 0.379216\\
029 & 5 & 9,383 & 9,383 & 0 & 0 & 0.000000\\
030 & 1 & 2,808,034 & 2,807,891 & 5,519,360 & 5,519,360 & 0.259868\\
030 & 2 & 1,400,031 & 1,400,031 & 2,354,176 & 2,354,176 & 0.073986\\
030 & 3 & 931,803 & 931,803 & 4,257,792 & 4,257,792 & 0.098845\\
030 & 4 & 696,636 & 696,606 & 5,032,960 & 5,032,960 & 0.649545\\
030 & 5 & 558,845 & 558,512 & 6,318,080 & 6,318,080 & 0.709768\\
031 & 1 & 1,403,390 & 1,403,039 & 3,281,920 & 3,281,920 & 0.183790\\
031 & 2 & 708,525 & 707,774 & 5,152,768 & 5,152,768 & 0.386861\\
031 & 3 & 471,415 & 471,340 & 3,615,744 & 3,615,744 & 0.531457\\
031 & 4 & 355,075 & 355,023 & 5,144,064 & 5,144,064 & 0.582847\\
031 & 5 & 287,749 & 284,116 & 5,666,304 & 5,666,304 & 0.627811\\
032 & 1 & 34,065 & 34,065 & 0 & 0 & 0.000000\\
032 & 2 & 17,286 & 17,286 & 6,216 & 6,216 & 0.021782\\
032 & 3 & 11,775 & 11,775 & 12,288 & 12,288 & 0.042162\\
032 & 4 & 9,481 & 9,481 & 12,300 & 12,300 & 0.028108\\
032 & 5 & 8,679 & 8,396 & 59,520 & 59,520 & 0.175604\\
033 & 1 & 1,769,406 & 1,769,212 & 4,737,024 & 4,737,024 & 0.198374\\
033 & 2 & 884,553 & 883,778 & 5,077,536 & 5,077,536 & 0.408071\\
033 & 3 & 587,457 & 587,082 & 7,426,184 & 7,426,184 & 0.315553\\
033 & 4 & 443,314 & 442,630 & 7,214,688 & 7,214,688 & 0.432788\\
033 & 5 & 350,567 & 350,553 & 8,951,424 & 8,951,424 & 0.592653\\
034 & 1 & 824,227 & 823,980 & 427,008 & 427,008 & 0.096840\\
034 & 2 & 391,716 & 391,580 & 2,076,168 & 2,076,168 & 0.278494\\
034 & 3 & 266,018 & 265,966 & 2,320,384 & 2,320,384 & 0.445850\\
034 & 4 & 196,606 & 196,515 & 2,861,568 & 2,861,568 & 0.454985\\
034 & 5 & 158,585 & 157,798 & 3,510,272 & 3,510,272 & 0.501897\\
035 & 1 & 165,528 & 165,528 & 0 & 0 & 0.000000\\
035 & 2 & 92,022 & 92,022 & 189,696 & 189,696 & 0.000000\\
035 & 3 & 79,468 & 79,468 & 584,908 & 584,908 & 0.154835\\
035 & 4 & 68,907 & 68,430 & 1,097,894 & 1,097,894 & 0.272892\\
035 & 5 & 65,452 & 65,383 & 718,326 & 718,326 & 0.172523\\
036 & 1 & 931,510 & 931,510 & 0 & 0 & 0.000000\\
036 & 2 & 466,068 & 466,068 & 1,152 & 1,152 & 0.000000\\
036 & 3 & 311,345 & 311,345 & 2,304 & 2,304 & 0.000000\\
036 & 4 & 233,584 & 233,584 & 3,456 & 3,456 & 0.000000\\
036 & 5 & 187,182 & 187,182 & 4,608 & 4,608 & 0.000000\\
037 & 1 & 210,001 & 210,001 & 0 & 0 & 0.000000\\
037 & 2 & 105,374 & 105,374 & 204,800 & 204,800 & 0.000000\\
037 & 3 & 71,271 & 71,193 & 1,390,592 & 1,390,592 & 0.525215\\
037 & 4 & 53,320 & 53,320 & 614,400 & 614,400 & 0.000000\\
037 & 5 & 43,013 & 43,013 & 901,120 & 901,120 & 0.378571\\
038 & 1 & 1,353,382 & 1,352,555 & 3,683,328 & 3,683,328 & 0.451944\\
038 & 2 & 674,500 & 674,382 & 6,790,656 & 6,790,656 & 0.278831\\
038 & 3 & 452,323 & 451,791 & 5,426,692 & 5,426,692 & 0.450613\\
038 & 4 & 338,878 & 338,219 & 5,772,288 & 5,772,288 & 0.625507\\
038 & 5 & 273,839 & 272,000 & 7,279,108 & 7,279,108 & 0.639257\\
039 & 1 & 1,425,165 & 1,424,699 & 8,220,672 & 8,220,672 & 0.022090\\
039 & 2 & 701,268 & 698,189 & 9,957,376 & 9,957,376 & 0.414641\\
039 & 3 & 474,891 & 468,365 & 11,755,520 & 11,755,520 & 0.533522\\
039 & 4 & 352,557 & 349,027 & 12,656,640 & 12,656,640 & 0.602975\\
039 & 5 & 336,649 & 279,449 & 14,704,640 & 14,704,640 & 0.612853\\
040 & 1 & 244,619 & 244,619 & 244,408 & 244,408 & 0.000000\\
040 & 2 & 130,452 & 130,452 & 1,402,390 & 1,402,390 & 0.030110\\
040 & 3 & 101,669 & 101,669 & 1,605,656 & 1,605,656 & 0.125432\\
040 & 4 & 86,134 & 85,446 & 1,720,946 & 1,720,946 & 0.320978\\
040 & 5 & 78,715 & 75,955 & 1,592,310 & 1,592,310 & 0.350297\\
041 & 1 & 5,986,657 & 5,984,182 & 22,643,712 & 22,643,712 & 0.138945\\
041 & 2 & 2,985,413 & 2,984,306 & 31,470,608 & 31,470,608 & 0.159201\\
041 & 3 & 1,991,784 & 1,988,631 & 35,111,940 & 35,111,940 & 0.270615\\
041 & 4 & 1,494,537 & 1,490,581 & 36,428,816 & 36,428,816 & 0.306092\\
041 & 5 & 1,196,918 & 1,194,456 & 40,669,188 & 40,669,188 & 0.324990\\
042 & 1 & 1,654,414 & 1,651,239 & 6,491,136 & 6,491,136 & 0.300077\\
042 & 2 & 833,960 & 831,867 & 7,068,960 & 7,068,960 & 0.359559\\
042 & 3 & 555,088 & 552,280 & 7,577,152 & 7,577,152 & 0.396720\\
042 & 4 & 423,126 & 420,768 & 11,353,676 & 11,353,676 & 0.475553\\
042 & 5 & 342,520 & 339,457 & 10,187,328 & 10,187,328 & 0.493709\\
043 & 1 & 656,875 & 656,850 & 2,483,184 & 2,483,184 & 0.032215\\
043 & 2 & 328,008 & 328,008 & 506,880 & 506,880 & 0.166450\\
043 & 3 & 246,360 & 246,359 & 674,748 & 674,748 & 0.257460\\
043 & 4 & 221,356 & 221,351 & 868,500 & 868,500 & 0.281457\\
043 & 5 & 217,586 & 206,007 & 6,002,246 & 6,002,246 & 0.436832\\
044 & 1 & 75,173 & 75,173 & 0 & 0 & 0.000000\\
044 & 2 & 43,931 & 43,882 & 930,400 & 930,400 & 0.001512\\
044 & 3 & 62,619 & 36,208 & 2,194,368 & 2,194,368 & 0.561360\\
044 & 4 & 62,619 & 36,208 & 2,194,368 & 2,194,368 & 0.561360\\
044 & 5 & 62,619 & 36,208 & 2,194,368 & 2,194,368 & 0.561360\\
045 & 1 & 125,614 & 125,614 & 0 & 0 & 0.000000\\
045 & 2 & 63,059 & 63,059 & 0 & 0 & 0.000000\\
045 & 3 & 42,228 & 42,228 & 20,480 & 20,480 & 0.014085\\
045 & 4 & 31,871 & 31,871 & 0 & 0 & 0.000000\\
045 & 5 & 25,770 & 25,770 & 40,960 & 40,960 & 0.027027\\
046 & 1 & 211,358 & 211,358 & 0 & 0 & 0.000000\\
046 & 2 & 106,042 & 106,042 & 930,400 & 930,400 & 0.000000\\
046 & 3 & 81,870 & 77,417 & 2,220,224 & 2,220,224 & 0.557003\\
046 & 4 & 88,823 & 62,397 & 2,905,888 & 2,905,888 & 0.501665\\
046 & 5 & 95,050 & 49,119 & 3,527,040 & 3,527,040 & 0.713206\\
047 & 1 & 303,591 & 303,591 & 0 & 0 & 0.000000\\
047 & 2 & 165,489 & 165,489 & 1,410,854 & 1,410,854 & 0.092853\\
047 & 3 & 126,353 & 126,167 & 2,029,640 & 2,029,640 & 0.104635\\
047 & 4 & 108,869 & 108,029 & 2,337,904 & 2,337,904 & 0.238281\\
047 & 5 & 105,845 & 102,266 & 2,927,124 & 2,927,124 & 0.229404\\
048 & 1 & 201,196 & 201,196 & 0 & 0 & 0.000000\\
048 & 2 & 135,478 & 135,478 & 757,062 & 757,062 & 0.107813\\
048 & 3 & 104,892 & 104,892 & 996,810 & 996,810 & 0.065659\\
048 & 4 & 104,892 & 104,892 & 996,810 & 996,810 & 0.065659\\
048 & 5 & 99,339 & 98,766 & 1,264,536 & 1,264,536 & 0.215307\\
049 & 1 & 181,926 & 181,926 & 278,036 & 278,036 & 0.106484\\
049 & 2 & 108,567 & 108,567 & 744,682 & 744,682 & 0.037405\\
049 & 3 & 83,467 & 83,117 & 1,008,040 & 1,008,040 & 0.146797\\
049 & 4 & 71,766 & 71,220 & 1,837,534 & 1,837,534 & 0.178255\\
049 & 5 & 68,665 & 67,575 & 2,011,460 & 2,011,460 & 0.204835\\
050 & 1 & 143,651 & 143,136 & 665,104 & 665,104 & 0.151713\\
050 & 2 & 78,745 & 78,745 & 485,524 & 485,524 & 0.073976\\
050 & 3 & 59,247 & 59,063 & 1,204,662 & 1,204,662 & 0.277186\\
050 & 4 & 53,438 & 51,949 & 1,540,648 & 1,540,648 & 0.273828\\
050 & 5 & 49,272 & 47,546 & 1,520,298 & 1,520,298 & 0.297621\\
051 & 1 & 607,628 & 607,628 & 0 & 0 & 0.000000\\
051 & 2 & 318,596 & 318,295 & 393,884 & 393,884 & 0.499788\\
051 & 3 & 227,412 & 227,412 & 952 & 952 & 0.000117\\
051 & 4 & 178,025 & 178,025 & 1,140 & 1,140 & 0.000234\\
051 & 5 & 169,451 & 167,637 & 493,392 & 493,392 & 0.555124\\
052 & 1 & 169,942 & 169,942 & 91,648 & 91,648 & 0.000000\\
052 & 2 & 80,154 & 80,005 & 1,119,386 & 1,119,386 & 0.174361\\
052 & 3 & 58,856 & 57,447 & 1,282,654 & 1,282,654 & 0.217224\\
052 & 4 & 44,668 & 44,460 & 632,750 & 632,750 & 0.269670\\
052 & 5 & 40,426 & 37,287 & 950,052 & 950,052 & 0.279195\\
053 & 1 & 1,367,418 & 1,367,418 & 65,098 & 65,098 & 0.000002\\
053 & 2 & 707,951 & 707,762 & 4,566,686 & 4,566,686 & 0.071320\\
053 & 3 & 497,588 & 485,409 & 7,234,968 & 7,234,968 & 0.292969\\
053 & 4 & 401,531 & 384,510 & 9,301,518 & 9,301,518 & 0.337458\\
053 & 5 & 388,730 & 361,811 & 12,007,080 & 12,007,080 & 0.339482\\
054 & 1 & 790,910 & 790,910 & 939,546 & 939,546 & 0.000010\\
054 & 2 & 441,988 & 441,987 & 4,365,002 & 4,365,002 & 0.029034\\
054 & 3 & 332,497 & 332,438 & 6,193,716 & 6,193,716 & 0.083756\\
054 & 4 & 286,675 & 286,390 & 7,240,730 & 7,240,730 & 0.092055\\
054 & 5 & 244,212 & 242,741 & 7,988,828 & 7,988,828 & 0.125910\\
055 & 1 & 103,144 & 102,631 & 1,429,632 & 1,429,632 & 0.120595\\
055 & 2 & 55,868 & 55,868 & 78,348 & 78,348 & 0.083391\\
055 & 3 & 37,611 & 37,611 & 115,712 & 115,712 & 0.156071\\
055 & 4 & 34,622 & 30,914 & 580,372 & 580,372 & 0.200066\\
055 & 5 & 23,198 & 23,160 & 173,568 & 173,568 & 0.217161\\
056 & 1 & 259,765 & 259,765 & 9,218 & 9,218 & 0.000009\\
056 & 2 & 163,006 & 163,006 & 721,074 & 721,074 & 0.074458\\
056 & 3 & 133,200 & 132,899 & 1,819,948 & 1,819,948 & 0.168137\\
056 & 4 & 121,296 & 121,253 & 2,402,690 & 2,402,690 & 0.184659\\
056 & 5 & 107,843 & 107,262 & 2,832,584 & 2,832,584 & 0.246831\\
057 & 1 & 58,573 & 58,573 & 73,728 & 73,728 & 0.063711\\
057 & 2 & 29,387 & 29,387 & 20,256 & 20,256 & 0.045547\\
057 & 3 & 19,648 & 19,648 & 392,978 & 392,978 & 0.466117\\
057 & 4 & 16,054 & 15,544 & 447,792 & 447,792 & 0.516481\\
057 & 5 & 12,956 & 12,956 & 216,640 & 216,640 & 0.035753\\
058 & 1 & 4,525,946 & 4,525,946 & 28,942,336 & 28,942,336 & 0.000000\\
058 & 2 & 2,260,738 & 2,247,911 & 33,566,720 & 33,566,720 & 0.285569\\
058 & 3 & 1,504,784 & 1,495,460 & 31,043,584 & 31,043,584 & 0.247171\\
058 & 4 & 1,138,267 & 1,118,279 & 31,543,296 & 31,543,296 & 0.487229\\
058 & 5 & 986,488 & 911,574 & 42,000,384 & 42,000,384 & 0.376295\\
059 & 1 & 738,041 & 738,041 & 0 & 0 & 0.000000\\
059 & 2 & 369,318 & 369,318 & 360,448 & 360,448 & 0.000000\\
059 & 3 & 250,622 & 250,289 & 1,363,968 & 1,363,968 & 0.523156\\
059 & 4 & 185,216 & 185,216 & 1,081,344 & 1,081,344 & 0.000000\\
059 & 5 & 150,386 & 150,246 & 5,713,920 & 5,713,920 & 0.322672\\
060 & 1 & 860,316 & 860,185 & 1,757,312 & 1,757,312 & 0.391465\\
060 & 2 & 436,355 & 436,329 & 2,357,504 & 2,357,504 & 0.423533\\
060 & 3 & 268,681 & 268,559 & 2,908,416 & 2,908,416 & 0.345030\\
060 & 4 & 208,430 & 207,019 & 3,093,632 & 3,093,632 & 0.412397\\
060 & 5 & 173,991 & 173,733 & 3,576,064 & 3,576,064 & 0.465576\\
061 & 1 & 77,191 & 77,191 & 0 & 0 & 0.000000\\
061 & 2 & 38,453 & 38,453 & 34,560 & 34,560 & 0.059731\\
061 & 3 & 27,905 & 27,905 & 33,792 & 33,792 & 0.058481\\
061 & 4 & 20,940 & 20,940 & 91,136 & 91,136 & 0.121291\\
061 & 5 & 17,354 & 16,981 & 104,960 & 104,960 & 0.124246\\
062 & 1 & 2,854,809 & 2,854,809 & 0 & 0 & 0.000000\\
062 & 2 & 1,428,964 & 1,428,964 & 75,264 & 75,264 & 0.000000\\
062 & 3 & 963,935 & 963,935 & 73,728 & 73,728 & 0.000000\\
062 & 4 & 750,419 & 750,199 & 22,798,848 & 22,798,848 & 0.000309\\
062 & 5 & 587,941 & 587,941 & 82,944 & 82,944 & 0.000000\\
063 & 1 & 1,070,397 & 1,070,397 & 0 & 0 & 0.000000\\
063 & 2 & 536,109 & 536,080 & 75,264 & 75,264 & 0.002730\\
063 & 3 & 358,787 & 358,787 & 55,296 & 55,296 & 0.000000\\
063 & 4 & 274,542 & 274,542 & 64,512 & 64,512 & 0.000000\\
063 & 5 & 221,161 & 221,161 & 67,584 & 67,584 & 0.000000\\
064 & 1 & 16,886 & 16,886 & 0 & 0 & 0.000000\\
064 & 2 & 12,557 & 12,556 & 46,214 & 46,214 & 0.157201\\
064 & 3 & 10,820 & 10,820 & 59,916 & 59,916 & 0.217036\\
064 & 4 & 9,630 & 9,629 & 94,446 & 94,446 & 0.282726\\
064 & 5 & 9,235 & 9,227 & 95,122 & 95,122 & 0.328570\\
065 & 1 & 45,516 & 45,516 & 0 & 0 & 0.000000\\
065 & 2 & 22,028 & 22,028 & 143,036 & 143,036 & 0.013824\\
065 & 3 & 14,713 & 14,513 & 85,056 & 85,056 & 0.037559\\
065 & 4 & 12,027 & 11,713 & 119,360 & 119,360 & 0.119939\\
065 & 5 & 10,578 & 10,267 & 74,160 & 74,160 & 0.140184\\
066 & 1 & 326,438 & 326,438 & 6,144 & 6,144 & 0.000000\\
066 & 2 & 193,524 & 193,440 & 1,350,136 & 1,350,136 & 0.117908\\
066 & 3 & 139,520 & 138,440 & 2,665,988 & 2,665,988 & 0.234074\\
066 & 4 & 119,686 & 117,281 & 2,438,828 & 2,438,828 & 0.316284\\
066 & 5 & 107,318 & 104,620 & 2,622,404 & 2,622,404 & 0.334707\\
067 & 1 & 61,063,003 & 61,063,003 & 34,926,592 & 34,926,592 & 0.008717\\
067 & 2 & 30,129,171 & 30,117,593 & 50,086,272 & 50,086,272 & 0.081504\\
067 & 3 & 19,982,761 & 19,951,926 & 58,923,776 & 58,923,776 & 0.156628\\
067 & 4 & 15,047,096 & 14,997,101 & 56,276,096 & 56,276,096 & 0.231720\\
067 & 5 & 12,074,623 & 12,061,334 & 71,656,960 & 71,656,960 & 0.199305\\
068 & 1 & 339,156 & 339,156 & 323,896 & 323,896 & 0.002819\\
068 & 2 & 191,661 & 191,254 & 2,000,138 & 2,000,138 & 0.068798\\
068 & 3 & 151,377 & 147,114 & 3,326,302 & 3,326,302 & 0.116865\\
068 & 4 & 133,321 & 127,752 & 2,789,758 & 2,789,758 & 0.203040\\
068 & 5 & 129,496 & 121,498 & 3,022,320 & 3,022,320 & 0.227078\\
069 & 1 & 21,899 & 21,899 & 0 & 0 & 0.000000\\
069 & 2 & 13,821 & 13,816 & 118,438 & 118,438 & 0.262749\\
069 & 3 & 10,068 & 10,008 & 253,194 & 253,194 & 0.312694\\
069 & 4 & 10,068 & 10,008 & 253,194 & 253,194 & 0.312694\\
069 & 5 & 8,472 & 7,886 & 320,114 & 320,114 & 0.252631\\
070 & 1 & 109,766 & 109,766 & 0 & 0 & 0.000000\\
070 & 2 & 51,604 & 51,503 & 168,468 & 168,468 & 0.232241\\
070 & 3 & 39,226 & 39,077 & 345,136 & 345,136 & 0.165974\\
070 & 4 & 30,029 & 29,793 & 346,636 & 346,636 & 0.344011\\
070 & 5 & 27,080 & 24,137 & 435,200 & 435,200 & 0.307778\\
071 & 1 & 16,084 & 16,084 & 0 & 0 & 0.000000\\
071 & 2 & 9,892 & 9,783 & 244,966 & 244,966 & 0.209102\\
071 & 3 & 8,660 & 8,430 & 227,976 & 227,976 & 0.299120\\
071 & 4 & 7,671 & 7,102 & 227,436 & 227,436 & 0.265096\\
071 & 5 & 7,004 & 5,619 & 272,274 & 272,274 & 0.290587\\
072 & 1 & 23,692,244 & 23,669,887 & 89,090,992 & 89,090,992 & 0.053645\\
072 & 2 & 11,821,325 & 11,807,422 & 104,426,272 & 104,426,272 & 0.219197\\
072 & 3 & 7,896,352 & 7,866,605 & 112,434,990 & 112,434,990 & 0.305969\\
072 & 4 & 5,969,952 & 5,924,972 & 122,580,790 & 122,580,790 & 0.372438\\
072 & 5 & 4,822,104 & 4,759,988 & 131,848,768 & 131,848,768 & 0.243522\\
073 & 1 & 8,221,623 & 8,221,623 & 188,416 & 188,416 & 0.000000\\
073 & 2 & 4,191,018 & 4,184,375 & 8,734,592 & 8,734,592 & 0.186503\\
073 & 3 & 2,793,228 & 2,759,896 & 7,844,608 & 7,844,608 & 0.621095\\
073 & 4 & 2,225,808 & 2,125,171 & 11,164,800 & 11,164,800 & 0.741808\\
073 & 5 & 1,804,127 & 1,738,206 & 20,287,488 & 20,287,488 & 0.479642\\
074 & 1 & 1,645,640 & 1,644,373 & 6,479,872 & 6,479,872 & 0.112620\\
074 & 2 & 813,707 & 813,107 & 5,988,360 & 5,988,360 & 0.326870\\
074 & 3 & 545,017 & 544,852 & 8,351,772 & 8,351,772 & 0.289935\\
074 & 4 & 407,804 & 405,974 & 7,229,452 & 7,229,452 & 0.315381\\
074 & 5 & 338,446 & 330,165 & 12,501,052 & 12,501,052 & 0.405933\\
075 & 1 & 1,358,701 & 1,358,701 & 0 & 0 & 0.000000\\
075 & 2 & 706,423 & 706,423 & 3,238,406 & 3,238,406 & 0.124760\\
075 & 3 & 581,136 & 579,216 & 11,876,108 & 11,876,108 & 0.235288\\
075 & 4 & 464,904 & 460,158 & 10,952,338 & 10,952,338 & 0.221981\\
075 & 5 & 387,016 & 369,043 & 12,400,568 & 12,400,568 & 0.270773\\
076 & 1 & 6,156,105 & 6,153,011 & 15,784,960 & 15,784,960 & 0.391238\\
076 & 2 & 3,085,509 & 3,082,459 & 26,401,792 & 26,401,792 & 0.391927\\
076 & 3 & 2,047,800 & 2,047,158 & 18,502,656 & 18,502,656 & 0.174656\\
076 & 4 & 1,529,160 & 1,529,141 & 20,473,856 & 20,473,856 & 0.146174\\
076 & 5 & 1,224,176 & 1,224,141 & 23,614,464 & 23,614,464 & 0.146676\\
077 & 1 & 488,026 & 488,026 & 0 & 0 & 0.000000\\
077 & 2 & 248,734 & 248,734 & 57,344 & 57,344 & 0.026082\\
077 & 3 & 169,489 & 169,489 & 114,688 & 114,688 & 0.025419\\
077 & 4 & 129,590 & 126,806 & 114,688 & 114,688 & 0.050838\\
077 & 5 & 115,617 & 115,617 & 172,032 & 172,032 & 0.074367\\
078 & 1 & 2,067,989 & 2,067,989 & 0 & 0 & 0.000000\\
078 & 2 & 1,024,634 & 1,024,634 & 2,148,480 & 2,148,480 & 0.000000\\
078 & 3 & 718,188 & 701,842 & 7,239,360 & 7,239,360 & 0.409308\\
078 & 4 & 521,160 & 517,424 & 6,973,824 & 6,973,824 & 0.498840\\
078 & 5 & 485,538 & 433,782 & 12,005,376 & 12,005,376 & 0.612358\\
079 & 1 & 30,856,884 & 30,856,884 & 32,772,096 & 32,772,096 & 0.000796\\
079 & 2 & 15,389,092 & 15,388,269 & 31,875,072 & 31,875,072 & 0.064450\\
079 & 3 & 10,117,508 & 10,113,306 & 38,875,136 & 38,875,136 & 0.144913\\
079 & 4 & 7,551,474 & 7,550,155 & 42,356,736 & 42,356,736 & 0.153231\\
079 & 5 & 6,063,040 & 6,062,064 & 33,882,112 & 33,882,112 & 0.467134\\
080 & 1 & 290,833 & 290,734 & 1,040,384 & 1,040,384 & 0.310153\\
080 & 2 & 145,406 & 145,338 & 1,785,856 & 1,785,856 & 0.378468\\
080 & 3 & 100,325 & 100,325 & 2,285,568 & 2,285,568 & 0.424242\\
080 & 4 & 80,344 & 76,782 & 1,708,032 & 1,708,032 & 0.289141\\
080 & 5 & 94,963 & 63,001 & 3,616,768 & 3,616,768 & 0.503110\\
081 & 1 & 151,795 & 151,795 & 0 & 0 & 0.000000\\
081 & 2 & 76,026 & 76,026 & 1,392,980 & 1,392,980 & 0.000000\\
081 & 3 & 51,943 & 51,943 & 0 & 0 & 0.000000\\
081 & 4 & 39,588 & 39,588 & 0 & 0 & 0.000000\\
081 & 5 & 30,943 & 30,943 & 0 & 0 & 0.000000\\
082 & 1 & 678,294 & 678,162 & 538,506 & 538,506 & 0.030927\\
082 & 2 & 348,139 & 348,070 & 1,623,014 & 1,623,014 & 0.099013\\
082 & 3 & 265,681 & 265,681 & 2,148,552 & 2,148,552 & 0.145167\\
082 & 4 & 227,802 & 226,842 & 3,139,858 & 3,139,858 & 0.349022\\
082 & 5 & 186,393 & 184,003 & 3,820,568 & 3,820,568 & 0.405185\\
083 & 1 & 841,782 & 841,782 & 131,072 & 131,072 & 0.000000\\
083 & 2 & 411,142 & 410,764 & 5,288,544 & 5,288,544 & 0.078177\\
083 & 3 & 282,133 & 281,995 & 3,277,760 & 3,277,760 & 0.227133\\
083 & 4 & 211,087 & 211,087 & 2,791,200 & 2,791,200 & 0.701874\\
083 & 5 & 203,523 & 177,330 & 5,802,368 & 5,802,368 & 0.575494\\
084 & 1 & 1,371,156 & 1,371,156 & 9,676,800 & 9,676,800 & 0.000000\\
084 & 2 & 686,336 & 686,336 & 9,925,632 & 9,925,632 & 0.000000\\
084 & 3 & 458,172 & 458,172 & 9,179,136 & 9,179,136 & 0.000000\\
084 & 4 & 353,412 & 353,412 & 737,280 & 737,280 & 0.000000\\
084 & 5 & 285,686 & 285,686 & 322,560 & 322,560 & 0.000000\\
085 & 1 & 3,097,342 & 3,097,342 & 3,983,714 & 3,983,714 & 0.000007\\
085 & 2 & 1,573,135 & 1,573,135 & 7,725,538 & 7,725,538 & 0.083973\\
085 & 3 & 1,193,752 & 1,193,708 & 10,444,556 & 10,444,556 & 0.140153\\
085 & 4 & 996,579 & 986,045 & 29,437,680 & 29,437,680 & 0.202026\\
085 & 5 & 897,610 & 862,197 & 37,590,868 & 37,590,868 & 0.289567\\
086 & 1 & 108,508 & 108,508 & 72,390 & 72,390 & 0.000017\\
086 & 2 & 61,340 & 61,165 & 1,253,446 & 1,253,446 & 0.094194\\
086 & 3 & 53,702 & 51,974 & 1,759,116 & 1,759,116 & 0.188586\\
086 & 4 & 48,576 & 44,949 & 2,088,146 & 2,088,146 & 0.331584\\
086 & 5 & 48,104 & 44,818 & 2,088,146 & 2,088,146 & 0.331333\\
087 & 1 & 3,952,857 & 3,948,565 & 8,317,770 & 8,317,770 & 0.180391\\
087 & 2 & 1,982,575 & 1,979,975 & 9,768,960 & 9,768,960 & 0.248793\\
087 & 3 & 1,346,667 & 1,345,661 & 10,874,522 & 10,874,522 & 0.244516\\
087 & 4 & 1,045,358 & 1,045,358 & 9,695,232 & 9,695,232 & 0.174712\\
087 & 5 & 832,158 & 829,663 & 12,487,680 & 12,487,680 & 0.299667\\
088 & 1 & 231,549 & 231,497 & 391,168 & 391,168 & 0.006049\\
088 & 2 & 132,283 & 131,063 & 1,343,442 & 1,343,442 & 0.081882\\
088 & 3 & 101,562 & 99,048 & 1,895,308 & 1,895,308 & 0.347745\\
088 & 4 & 88,613 & 85,617 & 1,807,122 & 1,807,122 & 0.125486\\
088 & 5 & 78,860 & 74,720 & 2,089,020 & 2,089,020 & 0.199440\\
089 & 1 & 593,264 & 593,264 & 0 & 0 & 0.000000\\
089 & 2 & 295,856 & 295,856 & 238,720 & 238,720 & 0.000000\\
089 & 3 & 203,941 & 203,894 & 1,806,592 & 1,806,592 & 0.165208\\
089 & 4 & 148,856 & 148,856 & 716,160 & 716,160 & 0.000000\\
089 & 5 & 120,170 & 120,170 & 954,880 & 954,880 & 0.000000\\
090 & 1 & 843,811 & 843,811 & 0 & 0 & 0.000000\\
090 & 2 & 453,180 & 451,406 & 3,341,952 & 3,341,952 & 0.125258\\
090 & 3 & 332,880 & 318,468 & 5,760,384 & 5,760,384 & 0.223952\\
090 & 4 & 288,180 & 255,234 & 6,606,720 & 6,606,720 & 0.537441\\
090 & 5 & 355,314 & 196,502 & 13,899,264 & 13,899,264 & 0.801573\\
091 & 1 & 9,113,927 & 9,112,016 & 36,733,952 & 36,733,952 & 0.096910\\
091 & 2 & 4,529,778 & 4,527,598 & 49,288,192 & 49,288,192 & 0.337376\\
091 & 3 & 3,039,758 & 3,032,804 & 60,819,456 & 60,819,456 & 0.459209\\
091 & 4 & 2,297,361 & 2,280,824 & 65,764,500 & 65,764,500 & 0.306464\\
091 & 5 & 1,846,514 & 1,828,683 & 43,944,448 & 43,944,448 & 0.558036\\
092 & 1 & 4,010,076 & 4,010,076 & 5,476,352 & 5,476,352 & 0.000000\\
092 & 2 & 2,009,876 & 2,009,860 & 6,070,880 & 6,070,880 & 0.044091\\
092 & 3 & 1,347,278 & 1,339,219 & 7,521,472 & 7,521,472 & 0.270440\\
092 & 4 & 1,017,001 & 996,428 & 6,264,608 & 6,264,608 & 0.280934\\
092 & 5 & 779,445 & 773,666 & 6,433,152 & 6,433,152 & 0.515603\\
093 & 1 & 74,205 & 74,205 & 0 & 0 & 0.000000\\
093 & 2 & 37,827 & 37,827 & 8,192 & 8,192 & 0.000000\\
093 & 3 & 25,103 & 25,103 & 16,384 & 16,384 & 0.000000\\
093 & 4 & 19,779 & 19,743 & 24,576 & 24,576 & 0.030822\\
093 & 5 & 15,841 & 15,595 & 36,864 & 36,864 & 0.040404\\
094 & 1 & 86,556 & 86,556 & 0 & 0 & 0.000000\\
094 & 2 & 39,937 & 38,992 & 115,200 & 115,200 & 0.087105\\
094 & 3 & 27,651 & 27,531 & 211,968 & 211,968 & 0.161851\\
094 & 4 & 21,777 & 21,738 & 318,816 & 318,816 & 0.123562\\
094 & 5 & 18,953 & 18,035 & 332,064 & 332,064 & 0.258464\\
095 & 1 & 1,144,413 & 1,144,413 & 16,809,984 & 16,809,984 & 0.000000\\
095 & 2 & 573,690 & 573,690 & 16,588,800 & 16,588,800 & 0.000000\\
095 & 3 & 433,227 & 433,227 & 2,433,024 & 2,433,024 & 0.000000\\
095 & 4 & 335,913 & 335,913 & 147,840 & 147,840 & 0.000000\\
095 & 5 & 269,789 & 269,789 & 147,456 & 147,456 & 0.000000\\
096 & 1 & 73,473 & 73,473 & 22,584 & 22,584 & 0.000000\\
096 & 2 & 36,559 & 36,346 & 196,484 & 196,484 & 0.340914\\
096 & 3 & 26,694 & 26,191 & 290,508 & 290,508 & 0.460380\\
096 & 4 & 20,448 & 19,873 & 649,432 & 649,432 & 0.439262\\
096 & 5 & 18,907 & 18,630 & 442,444 & 442,444 & 0.475327\\
097 & 1 & 10,934,656 & 10,934,396 & 58,613,760 & 58,613,760 & 0.003654\\
097 & 2 & 5,400,820 & 5,398,228 & 58,441,728 & 58,441,728 & 0.074868\\
097 & 3 & 3,551,649 & 3,551,104 & 14,641,152 & 14,641,152 & 0.305868\\
097 & 4 & 2,597,292 & 2,596,536 & 17,580,032 & 17,580,032 & 0.301013\\
097 & 5 & 2,100,763 & 2,100,021 & 18,178,048 & 18,178,048 & 0.438414\\
098 & 1 & 166,777 & 166,777 & 372,804 & 372,804 & 0.013800\\
098 & 2 & 84,269 & 83,520 & 550,092 & 550,092 & 0.180050\\
098 & 3 & 57,758 & 57,758 & 707,860 & 707,860 & 0.196636\\
098 & 4 & 46,315 & 46,309 & 784,600 & 784,600 & 0.308274\\
098 & 5 & 39,900 & 39,669 & 830,752 & 830,752 & 0.262160\\
099 & 1 & 1,611,560 & 1,609,527 & 4,550,656 & 4,550,656 & 0.252665\\
099 & 2 & 809,724 & 806,783 & 7,327,872 & 7,327,872 & 0.420402\\
099 & 3 & 545,191 & 540,097 & 8,877,696 & 8,877,696 & 0.524879\\
099 & 4 & 406,980 & 403,923 & 9,291,392 & 9,291,392 & 0.566102\\
099 & 5 & 328,699 & 324,499 & 10,166,276 & 10,166,276 & 0.558698\\
100 & 1 & 122,765 & 122,765 & 278,128 & 278,128 & 0.014144\\
100 & 2 & 64,605 & 64,209 & 624,972 & 624,972 & 0.327489\\
100 & 3 & 47,629 & 46,369 & 570,616 & 570,616 & 0.490784\\
100 & 4 & 40,957 & 38,002 & 1,283,016 & 1,283,016 & 0.450185\\
100 & 5 & 35,783 & 32,716 & 777,764 & 777,764 & 0.492463\\
\end{longtable}
\endgroup

\AppendixPart{6}{算法参数的实验取值}\label{app:parameters}
以下取值对应正文所述的实际方法。时间参数以主机秒计，服务时间和交接阈值以模拟周期计。计数和权重均为无量纲量。题给硬件参数仍按正文执行模型设置。

\Needspace{8\baselineskip}
\begingroup\small
\setlength{\tabcolsep}{4pt}
\setlength{\LTleft}{0pt}\setlength{\LTright}{0pt}
\begin{longtable}{@{}p{36mm}p{57mm}p{59mm}@{}}
\caption{初始构造与搜索控制}\label{tab:parameters_search}\\
\toprule
参数 & 含义 & 实验取值\\
\midrule
\endfirsthead
\caption[]{初始构造与搜索控制（续）}\\
\toprule
参数 & 含义 & 实验取值\\
\midrule
\endhead
\bottomrule
\endfoot
$t_0,b_0$ & 最小服务周期、比例分母的最小字节量 & 1周期，1 B\\
$h_{\rm cut},\rho_{\rm cut}$ & 短序列阈值、切分比例下限 & 6，0.15\\
$\gamma_1,\gamma_2,\gamma_3$ & 三问的边界通信权重 & 2，1，1\\
$\mathcal G,\kappa_{\rm cut}$ & 构造粒度、切分次数系数 & $\{2,4,8\}$，2\\
$\lambda_{\rm in}$ & 初始核心映射的输入通信权重 & 0.25\\
$N_{\rm init},K_{\rm init}$ & 初始评价上限、每次补入候选上限 & 3，12\\
$\tau_{\rm init},\alpha_{\rm init}$ & 初始时间下限、预算比例 & 2 s，0.25\\
$\mathcal R_{\rm split}$ & 前缀切分比例集合 & $\{1/4,1/2,3/4\}$\\
$\delta_{\rm rel},K_{\rm rel}$ & 输出提前量阈值、候选上限 & 100周期，4\\
$\eta_{\rm exp},\nu_0,\nu_1$ & 探索强度、调用计数平滑量 & 0.012，2，1\\
$c_{\min},\lambda_{\rm fail}$ & 调用成本下限、失败惩罚权重 & 0.005 s，1\\
$w_{\rm hist},H_{\rm gap}$ & 近期统计窗口、优先再访间隔 & 8，8\\
$N_{\rm cover},K_{\rm cover}$ & 初次覆盖的调用数、候选数阈值 & 2，12\\
$N_{\rm warm}$ & FIFO来源启用前的成功调用阈值 & 4\\
$\tau_{\rm warm},\alpha_{\rm warm}$ & FIFO来源暂缓时限、预算比例 & 20 s，0.18\\
$\kappa_{\rm cost},\tau_{\rm margin}$ & 准入成本放大系数、时间余量 & 1.2，0.05 s\\
$\tau_{\max},N_{\max}$ & 自适应迭代总时间、评价次数上限 & 600 s，2048\\
$\epsilon_{\rm gap}$ & 相对当前周期的近界容差 & 0.0025\\
$\tau_{\rm stall},N_{\rm stall}$ & 连续无改善时间、评价次数 & 30 s，32\\
$N_{\rm stall}^{\rm src}$ & 各可用来源连续无改善次数 & 8\\
\end{longtable}\endgroup

\Needspace{8\baselineskip}
\begingroup\small
\setlength{\tabcolsep}{4pt}
\setlength{\LTleft}{0pt}\setlength{\LTright}{0pt}
\begin{longtable}{@{}p{36mm}p{57mm}p{59mm}@{}}
\caption{驻留重组与缓存候选}\label{tab:parameters_residency}\\
\toprule
参数 & 含义 & 实验取值\\
\midrule
\endfirsthead
\toprule
参数 & 含义 & 实验取值\\
\midrule
\endhead
\bottomrule
\endfoot
$c_{\rm comp},h_{\rm small}$ & 分量数下限、小分量大小上限 & 8，32\\
$\rho_{\rm small},\rho_{\rm share}$ & 小分量算子占比、共享输入比阈值 & 0.75，0.10\\
$K_{\rm tensor},K_{\rm hull},K_{\rm move}$ & 张量、闭包子图、迁移候选上限 & 20，10，20\\
$\mathcal B_{\rm owner}$ & 保持核心映射的带数尝试顺序 & $(16,32,64,8)$\\
$\mathcal A_{\rm batch}$ & 容量放大系数与阶段数尝试顺序 & $(1,2),(1.5,2),(2,2),(1,3),$\newline$(2.5,2),(1.5,3),(2,3),(2.5,3)$\\
$(\alpha_0,\alpha_1,\alpha_2)$ & FIFO请求排序的三项权重 & $(0.25,0.5,0.25)$\\
$K_{\rm req},K_{\rm look}$ & 优先请求数、相邻子图考察上限 & 12，32\\
$\mathcal D_{\rm use}$ & 使用算子所在子图的移动距离 & $\{1,2,4,8\}$\\
$\mathcal D_{\rm fill},\mathcal D_{\rm evict}$ & 填充、淘汰相关子图的移动距离 & $\{1,2,4\}$，$\{1,2\}$\\
$d_{\rm joint}$ & 联合移动的子图位移 & 1\\
$(K^{\rm F}_1,\ldots,K^{\rm F}_5)$ & 五步生成过程的累计候选上限 & $(12,18,22,24,32)$\\
\end{longtable}\endgroup

\Needspace{8\baselineskip}
\begingroup\small
\setlength{\tabcolsep}{4pt}
\setlength{\LTleft}{0pt}\setlength{\LTright}{0pt}
\begin{longtable}{@{}p{36mm}p{57mm}p{59mm}@{}}
\caption{第三问多粒度候选优选}\label{tab:parameters_portfolio}\\
\toprule
参数 & 含义 & 实验取值\\
\midrule
\endfirsthead
\toprule
参数 & 含义 & 实验取值\\
\midrule
\endhead
\bottomrule
\endfoot
$\alpha_{\rm band}$ & 当前深度层带负载权重 & 0.7\\
$\eta_{\rm cache}$ & 通信排序中的驻留收益权重 & 0.5\\
$\mathcal B_{\rm plain}$ & 普通负载映射的目标带数 & $\{2,4,16,32,64\}$\\
$\mathcal A_{\rm comm}$ & 通信感知映射的带数、权重组合 & $\{4,8,16,32,64\}\times\{0.25\}$\newline$\cup\{(8,0.75)\}$\\
$\mu_{\rm bal}$ & 完整分量映射的总负载权重 & 0.15\\
$\varepsilon$ & 候选优选的相对周期改善门槛 & $10^{-3}$\\
\end{longtable}\endgroup


\AppendixEnd
